\chapter[Marco teórico]{Marco teórico}
\label{cp:marco}
% Extensión estimada: sin límite (guía, §5).
% Estructura (revisión del autor, 2026-10-04): 2.1 fundamentos del aprendizaje profundo, 2.2 transformers y
% atención, 2.3 clasificación visual, 2.4 modelos de lenguaje, 2.5 modelos visión--lenguaje, 2.6 incertidumbre y
% predicción selectiva, 2.7 generación de reportes, 2.8 evaluación, 2.9 sesgo, robustez y generalización,
% 2.10 estado del arte y posicionamiento. Solo se presentan los conceptos que el trabajo usa (tutor, Audio 7), en
% orden cronológico y con una transición entre secciones (Audios 7 y 8).
% Fuentes: Docs/marco_teorico/Documento_LLMs_VLLMs.tex, Docs/estado_del_arte/resumen_papers.tex (fichas de
% lectura de ~40 trabajos; no es fuente de cifras) y Docs/tabla_resultados_sota.

En este capítulo se introducen los fundamentos teóricos y los conceptos necesarios para contextualizar y
comprender el desarrollo del presente trabajo. Asimismo, se realiza una revisión de los trabajos relacionados
más relevantes, con el propósito de situar la propuesta en el contexto de las investigaciones existentes.

La organización del capítulo sigue una secuencia que permite introducir progresivamente los conceptos
fundamentales, comenzando por los principios del aprendizaje en redes neuronales, para luego abordar las
arquitecturas desarrolladas para el procesamiento de imágenes y texto y, posteriormente, los modelos que
integran ambas modalidades. A partir de estos fundamentos, se presentan aspectos relacionados con la
incertidumbre de las predicciones, la generación de reportes clínicos y los métodos de evaluación de los
modelos, incluyendo sus alcances y limitaciones. Finalmente, se revisa el estado del arte y se sitúa en él el
presente trabajo.

\section{Fundamentos del aprendizaje profundo}
\label{sec:marco_fundamentos}

El aprendizaje automático estudia algoritmos que mejoran su desempeño en una tarea a medida que acumulan
experiencia, que en la práctica son datos. En lugar de programar la solución como un conjunto de reglas, se
elige un modelo con parámetros ajustables y se los fija a partir de ejemplos, de modo que el modelo resuelva la
tarea también sobre datos que no vio~\cite{goodfellow2016deep}. En el aprendizaje \textbf{supervisado}, cada
ejemplo es un par formado por una entrada, como una imagen, un texto o un conjunto de mediciones, y la salida
esperada, su etiqueta, y el modelo aprende a predecir la etiqueta de entradas nuevas. En el \textbf{no
supervisado} los ejemplos no tienen etiqueta y se busca la estructura de los propios datos, como grupos de
casos parecidos~\cite{goodfellow2016deep}. El \textbf{autosupervisado} obtiene la etiqueta de los mismos datos,
por ejemplo la palabra que sigue en un texto, y es como se preentrenan los modelos de lenguaje
(Sección~\ref{sec:marco_llm_grandes}).

Durante décadas, el desempeño de estos métodos dependió de las características con que se describe cada
ejemplo. Un especialista las diseñaba a mano, como medidas de color o textura para una imagen o recuentos de
palabras para un texto, y el algoritmo aprendía solo a decidir a partir de ellas. El aprendizaje profundo
aprende también esas características, directamente de los píxeles o de las palabras. Para eso encadena muchas
transformaciones simples, de modo que las primeras capturan rasgos elementales y las últimas, conceptos más
abstractos; en una imagen, de bordes a contornos y de contornos a partes de objetos. El enfoque se impuso a
partir de la década de 2010 en el reconocimiento de imágenes y de voz y en el procesamiento del lenguaje,
cuando hubo suficientes datos y cómputo para entrenar esos modelos~\cite{goodfellow2016deep}, y es la base de
todos los modelos que usa este trabajo. Las transformaciones que encadena son las capas de una red neuronal.

\subsection{Redes neuronales y aprendizaje supervisado}
\label{sec:marco_redes}

Una red neuronal artificial es un modelo formado por unidades de cálculo simples, las neuronas, conectadas
entre sí y organizadas en capas. Cada neurona calcula una suma ponderada de sus entradas, le suma un sesgo y
aplica al resultado una función no lineal, llamada función de activación. Una capa reúne varias neuronas que reciben la misma entrada,
y un perceptrón multicapa apila capas de modo que la salida de cada una es la entrada de la siguiente. La capa
$l$ transforma el vector $\mathbf{h}^{(l-1)}$ que recibe en
\begin{equation}
    \mathbf{h}^{(l)} = \phi\!\left(W^{(l)} \mathbf{h}^{(l-1)} + \mathbf{b}^{(l)}\right),
    \label{eq:marco_capa}
\end{equation}
donde $W^{(l)}$ es la matriz de pesos, $\mathbf{b}^{(l)}$ el vector de sesgos y $\phi$ la función de activación,
aplicada por elemento. Sin la activación, la red entera se reduciría a una sola transformación lineal. Los pesos
y los sesgos de todas las capas son los parámetros de la red, agrupados en $\theta$, y una red de muchas capas
se llama profunda~\cite{goodfellow2016deep}.

Entrenar la red es elegir los valores de $\theta$. Para eso se define una función de pérdida $\mathcal{L}$, que
cuantifica cuánto se equivoca la red en un ejemplo, y se buscan los parámetros que minimizan la pérdida promedio
sobre el conjunto de entrenamiento. El descenso del gradiente lo hace de forma iterativa. El gradiente de la
pérdida respecto de los parámetros, $\nabla_{\theta} \mathcal{L}$, indica en qué dirección crece la pérdida, y
en cada paso los parámetros se desplazan en la dirección opuesta una distancia que fija la tasa de aprendizaje
$\eta$. La regla de actualización es entonces
\begin{equation}
    \theta \leftarrow \theta - \eta \, \nabla_{\theta} \mathcal{L}.
    \label{eq:marco_gradiente}
\end{equation}
Cada paso estima el gradiente sobre un lote de ejemplos y no sobre el conjunto completo, y una época es una
pasada por todos los lotes~\cite{goodfellow2016deep}.

El gradiente de todos los parámetros se obtiene con la retropropagación, que recorre la red desde la pérdida
hacia la entrada aplicando la regla de la cadena (Figura~\ref{fig:marco_perceptron}), con un costo del mismo orden que el de calcular la
salida~\cite{rumelhart1986learning}. Sobre la regla de la Ecuación~\ref{eq:marco_gradiente} se construyen
optimizadores como Adam, que adapta el paso de cada parámetro con dos promedios móviles, uno del gradiente y
otro de su cuadrado~\cite{kingma2015adam}. Esos promedios son memoria adicional por cada parámetro que se
entrena, un costo que pesa al adaptar modelos grandes (Sección~\ref{sec:marco_llm_lora}). La tasa de
aprendizaje suele crecer de forma lineal en las primeras épocas, un calentamiento (\emph{warmup}) que evita
pasos grandes mientras la red cambia rápido~\cite{goyal2017accurate}, y decaer después hasta cerca de cero
siguiendo medio ciclo de coseno~\cite{loshchilov2017sgdr}. Una vez entrenada, la red predice la salida de una
entrada nueva con una pasada hacia adelante y los parámetros fijos, lo que se llama inferencia.

% Figura: perceptrón multicapa, pasada hacia adelante y retropropagación.
% Se usa en sec:marco_fundamentos, subsección de entrenamiento.

\begin{figure}[H]
    \centering
    \begin{tikzpicture}[x=1mm, y=1mm,
        neurona/.style = {circle, draw=figtrazo, minimum size=6mm, inner sep=0pt, font=\footnotesize},
    ]

        % Capas.
        \foreach \i in {1, 2, 3} {
            \node[neurona, fill=figazul!15, draw=figazul!80!black] (x\i) at (0, 16 - \i*8) {$x_{\i}$};
        }
        \foreach \i in {1, 2, 3, 4} {
            \node[neurona, fill=fignaranja!18, draw=fignaranja!70!black] (h\i) at (25, 20 - \i*8) {};
        }
        \foreach \i in {1, 2, 3} {
            \node[neurona, fill=figverde!12, draw=figverde!70!black] (z\i) at (50, 16 - \i*8) {$z_{\i}$};
        }

        % Conexiones.
        \foreach \i in {1, 2, 3} { \foreach \j in {1, 2, 3, 4} {
            \draw[figtrazo!60, thin] (x\i) -- (h\j);
        } }
        \foreach \j in {1, 2, 3, 4} { \foreach \k in {1, 2, 3} {
            \draw[figtrazo!60, thin] (h\j) -- (z\k);
        } }

        % Softmax y pérdida.
        \node[figtarea, text width=16mm, minimum height=8mm] (sm) at (72, 0) {softmax};
        \node[figsalida, text width=20mm, minimum height=8mm] (loss) at (100, 0)
            {pérdida\\$\mathcal{L} = -\log p_y$};
        \draw[figflecha] (56, 0) -- (sm.west);
        \draw[figflecha] (sm.east) -- node[figetiqueta, above] {$p_c$} (loss.west);

        % Rótulos de capas.
        \node[font=\footnotesize, text=figtrazo] at (0, -22) {entrada};
        \node[font=\footnotesize, text=figtrazo, align=center] at (25, -22) {capa oculta\\$\phi(W\mathbf{x}+\mathbf{b})$};
        \node[font=\footnotesize, text=figtrazo] at (50, -22) {\emph{logits}};

        % Pasadas.
        \draw[figflecha, figazul!80!black, thick] (0, 24) -- node[figetiqueta, above, text=figazul!80!black]
            {pasada hacia adelante: se calcula la salida} (100, 24);
        \draw[-{Latex[length=2mm, width=1.6mm]}, thick, densely dashed, figbermellon]
            (100, -32) -- node[figetiqueta, below, text=figbermellon]
            {retropropagación: se calcula $\nabla_\theta \mathcal{L}$ capa por capa} (0, -32);

    \end{tikzpicture}
    \caption{Perceptrón multicapa con una capa oculta: pasada hacia adelante y retropropagación.}
    \label{fig:marco_perceptron}
\end{figure}


\subsection{Clasificación y funciones de pérdida}
\label{sec:marco_softmax}

Un problema de clasificación consiste en asignar cada entrada a una clase de un conjunto finito. En el diagnóstico de lesiones de piel, las clases son los tipos de lesión, como el
melanoma o el nevo, y el modelo recibe la imagen y debe indicar a cuál pertenece. Para un problema de $C$
clases, la última capa de la red tiene una salida por clase y devuelve un vector de $C$ puntajes reales sin
restricción, los \emph{logits} $\mathbf{z}$. La función \emph{softmax} los convierte en una distribución de
probabilidad,
\begin{equation}
    p_c = \frac{\exp(z_c)}{\sum_{j=1}^{C} \exp(z_j)}, \qquad c = 1, \dots, C,
    \label{eq:marco_softmax}
\end{equation}
con valores positivos que suman uno. La clase predicha es la de mayor probabilidad, y esa probabilidad se lee
como la confianza del modelo en su predicción, una lectura que solo vale si el modelo está calibrado
(Sección~\ref{sec:marco_senales_confianza}). Los modelos de lenguaje usan la misma operación para elegir cada
palabra entre las de su vocabulario (Sección~\ref{sec:marco_llm_grandes}).

La pérdida de clasificación es la entropía cruzada. Compara la distribución que predice el modelo,
$\mathbf{p}$, con la de la etiqueta, $\mathbf{q}$, que pone toda la probabilidad en la clase correcta $y$, y
como $q_c$ vale cero para las demás clases, se reduce a menos el logaritmo de la probabilidad asignada a la
correcta,
\begin{equation}
    \mathcal{L} = -\sum_{c=1}^{C} q_c \log p_c = -\log p_y.
    \label{eq:marco_entropia}
\end{equation}
Vale cero cuando toda la probabilidad está en la clase correcta y crece sin límite a medida que esa
probabilidad se acerca a cero, de modo que un error cometido con mucha confianza pesa más que uno cometido con
duda. Minimizarla equivale a maximizar la probabilidad que el modelo asigna a las respuestas
correctas~\cite{goodfellow2016deep}. Combinada con la softmax, su gradiente respecto de cada logit $z_c$ es
$p_c - q_c$: el entrenamiento sube el logit de la clase correcta y baja el de las demás en proporción a la
probabilidad que el modelo les había asignado. Los modelos de lenguaje usan la misma pérdida para cada token
(Sección~\ref{sec:marco_llm_grandes}), y las correcciones para el desbalance de clases
(Sección~\ref{sec:marco_desbalance}) son variantes de ella.

\subsection{Generalización y desafíos del aprendizaje}
\label{sec:marco_regularizacion}

Lo que interesa de un modelo es su desempeño sobre casos que no vio, y una red con suficientes parámetros puede
memorizar el conjunto de entrenamiento: la pérdida de entrenamiento sigue bajando mientras el error sobre datos
nuevos sube, fenómeno que se llama \textbf{sobreajuste}~\cite{goodfellow2016deep}.

Ese error se estima con ejemplos que el modelo no usó para ajustarse, y por eso los datos se dividen en tres
particiones disjuntas. La de entrenamiento ajusta los parámetros, la de validación elige entre configuraciones
e hiperparámetros, y la de test se usa una sola vez al final, porque una decisión tomada mirando el test vuelve
optimista su estimación~\cite{goodfellow2016deep}. Si imágenes de un mismo paciente o de una misma lesión caen
en el entrenamiento y en el test, el modelo puede acertar las segundas por haber visto las primeras y no porque
generalice. Esta fuga de información (\emph{data leakage}) hace la estimación más favorable que el desempeño
real, y se evita si todas las observaciones de un mismo caso caen en la misma partición.

La \textbf{regularización} reduce el sobreajuste restringiendo lo que la red puede aprender o ampliando la
variedad de lo que ve. El decaimiento de pesos (\emph{weight decay}) encoge los parámetros hacia cero en cada
paso, en una proporción $\lambda$, lo que con el descenso del gradiente equivale a sumar a la pérdida la
penalización $(\lambda/2)\sum_i \theta_i^2$~\cite{goodfellow2016deep}. Con Adam esa equivalencia no vale, y
AdamW aplica el decaimiento por separado de la actualización del optimizador~\cite{loshchilov2019decoupled}. La
detención temprana (\emph{early stopping}) conserva los parámetros de la época con mejor desempeño en
validación~\cite{goodfellow2016deep}. El suavizado de etiquetas (\emph{label smoothing}) le quita a la clase
correcta de $\mathbf{q}$ una fracción $\varepsilon$ de la probabilidad y la reparte entre todas las clases,
con lo que el modelo deja de empujar $p_y$ hacia uno y suele quedar mejor calibrado
(Sección~\ref{sec:marco_senales_confianza})~\cite{szegedy2016rethinking,muller2019does}.

El aumento de datos aplica a cada imagen, cada vez que entra en un lote, una transformación aleatoria que
preserva su clase, como una rotación, un recorte o un cambio de brillo. RandAugment aplica $N$ operaciones
sorteadas de una lista fija con una magnitud común $M$~\cite{cubuk2020randaugment}, y el borrado aleatorio
(\emph{random erasing}) tapa un rectángulo al azar para simular oclusiones~\cite{zhong2020random}. Cuando el color distingue las clases, como
en dermatoscopía~\cite{kittler2016standardization}, una transformación de color fuerte puede borrar la
información que las define~\cite{shorten2019survey}. El aumento se usa también al predecir: el \gls{tta}
promedia las probabilidades que el modelo asigna a la imagen original y a algunas versiones transformadas de
ella, como su reflejo horizontal~\cite{krizhevsky2012imagenet}.

Otro desafío es el desbalance de clases\label{sec:marco_desbalance}. Cuando unas clases tienen muchos más
ejemplos que otras, minimizar la entropía cruzada promedio favorece a las frecuentes, porque sus ejemplos
dominan la suma de la pérdida, y el efecto crece con el grado de desbalance~\cite{buda2018systematic}. Sean
$n$ la cantidad total de ejemplos de entrenamiento, $n_c$ la de la clase $c$ e $y$ la clase correcta del
ejemplo. Las correcciones más usadas modifican la entropía cruzada de la Ecuación~\ref{eq:marco_entropia} para
que las clases raras pesen más:
\begin{enumerate}
    \item la \textbf{pérdida ponderada}, $-w_y \log p_y$, multiplica la pérdida de cada ejemplo por un peso
    inversamente proporcional a la frecuencia de su clase, $w_y = n/(C\,n_y)$, de modo que cada clase aporta
    en total el mismo peso a la suma;
    \item la \textbf{pérdida focal}, $-(1-p_y)^{\gamma}\log p_y$, agrega un factor que se acerca a cero en los
    ejemplos que el modelo ya clasifica con confianza, con $p_y$ cerca de uno, y concentra así el
    entrenamiento en los difíciles. El parámetro $\gamma \ge 0$ regula la atenuación, y con $\gamma = 0$ se
    recupera la entropía cruzada~\cite{lin2017focal};
    \item el \textbf{\emph{balanced softmax}} multiplica en la softmax de la Ecuación~\ref{eq:marco_softmax}
    la exponencial de cada logit por la cantidad de ejemplos de su clase,
    $p_c = n_c \exp(z_c) / \sum_{j} n_j \exp(z_j)$, lo que equivale a sumar $\log n_c$ al logit $z_c$. Durante
    el entrenamiento, las clases frecuentes reciben así una ventaja que el modelo no necesita aprender en sus
    logits, y al predecir se usa la softmax común, sin ese término~\cite{ren2020balanced}.
\end{enumerate}
La alternativa es entrenar sin corrección y ajustar la regla de decisión, dividiendo la probabilidad de cada
clase por su proporción en el entrenamiento~\cite{buda2018systematic} o eligiendo en validación un umbral para
una clase rara. Ninguna de estas correcciones agrega información sobre las clases raras. Además, la exactitud
global oculta el desbalance, y por eso se usan métricas que promedian por clase
(Sección~\ref{sec:marco_eval_clasif}): la sensibilidad macro, que mide qué fracción de los casos de cada clase
se detecta, y el F1 macro, que además penaliza los falsos positivos.

\subsection{Modelos preentrenados y adaptación}
\label{sec:marco_transferencia}

Con pocos ejemplos etiquetados, como ocurre en la mayoría de los conjuntos de datos médicos, entrenar una red
desde cero favorece el sobreajuste. La transferencia de aprendizaje reutiliza lo aprendido en una tarea de
origen con muchos datos para mejorar el aprendizaje en la tarea de destino~\cite{pan2010survey}. En
reconocimiento visual, el preentrenamiento se hace habitualmente sobre ImageNet, un conjunto de referencia de
imágenes naturales etiquetadas~\cite{russakovsky2015imagenet,raghu2019transfusion}. Las primeras capas de una
red visual aprenden rasgos generales, parecidos a filtros de bordes y de color, que se reutilizan bien en otras
tareas, y las últimas, rasgos específicos de la tarea de origen~\cite{yosinski2014transferable}. En imágenes
médicas el beneficio puede ser menor al esperado, por la distancia con las imágenes
naturales~\cite{raghu2019transfusion}, y esa distancia motiva los modelos preentrenados sobre imágenes del
propio dominio, como DermLIP (Sección~\ref{sec:marco_vlm_clip}).

Para adaptar una red que clasifica se la divide en dos partes. El \emph{backbone} convierte la imagen en un
vector de características, y la cabeza de clasificación (\emph{head}) es la capa final, de la forma de la
Ecuación~\ref{eq:marco_capa}, que convierte ese vector en los logits. La cabeza se reemplaza por una nueva, con
una salida por cada clase de la tarea de destino. La \textbf{extracción de características} congela el
\emph{backbone}, es decir, mantiene fijos sus parámetros, y entrena solo la cabeza. El \textbf{ajuste fino}
(\emph{fine-tuning}) actualiza también el \emph{backbone}, completo o solo sus últimas capas. Entrenar más
parámetros da más libertad para adaptarse a la tarea nueva, pero con pocos datos aumenta el riesgo de
sobreajuste y de olvido catastrófico (\emph{catastrophic forgetting}), es decir, que el modelo pierda parte de
lo que aprendió en el preentrenamiento.

Entre esos dos extremos hay estrategias que adaptan el \emph{backbone} sin deformar lo que aprendió. La primera
entrena en dos etapas: primero solo la cabeza, con el \emph{backbone} congelado, y después la red completa. Si
el ajuste fino empieza con una cabeza inicializada al azar, los gradientes que llegan al \emph{backbone}
mientras la cabeza aprende deforman las características preentrenadas. Con la cabeza ya entrenada, la
secuencia rinde más que cada una de sus dos partes, sobre todo ante datos de una distribución distinta de la
del entrenamiento~\cite{kumar2022finetuning}. La segunda asigna a cada capa su propia tasa de aprendizaje,
decreciente hacia la entrada, porque las capas capturan información de distinto tipo~\cite{howard2018universal}.
Con $L$ capas numeradas desde la entrada y un factor $\xi < 1$, la capa $l$ se entrena con la tasa
\begin{equation}
    \eta_l = \eta \, \xi^{\,L-l}, \qquad l = 1, \dots, L,
    \label{eq:marco_llrd}
\end{equation}
de modo que la cabeza recibe la tasa completa y las primeras capas, que guardan los rasgos más generales,
cambian menos. La técnica se conoce como decaimiento de la tasa de aprendizaje por capa (\emph{layer-wise
learning rate decay}).

Hasta aquí, la red trata su entrada como un vector sin estructura. Las arquitecturas que hicieron práctico el
aprendizaje profundo incorporaron la estructura de cada tipo de dato, y la primera en hacerlo con las imágenes
fue la red convolucional.

\section{Arquitecturas para clasificación de imágenes}
\label{sec:marco_vision}

La red neuronal de la Sección~\ref{sec:marco_fundamentos} recibe su entrada como un vector y procesa cada valor
sin tener en cuenta dónde está ubicado. En una imagen, sin embargo, la ubicación es importante: los píxeles vecinos
suelen pertenecer a un mismo borde o a una misma región, y un mismo patrón, como el borde de una lesión, puede
aparecer en cualquier lugar. Las arquitecturas para clasificación de imágenes incorporan esta estructura de dos
maneras, la convolución y la atención. Los cuatro clasificadores que evalúa el trabajo pertenecen a esas dos
familias: DenseNet y ConvNeXt son redes convolucionales, y el \gls{vit} y el Swin Transformer se basan en la
atención. La sección presenta estas arquitecturas comenzando con las redes convolucionales y siguiendo con la
atención, el transformer y sus adaptaciones a imágenes.

\subsection{Redes convolucionales}
\label{sec:marco_cnn}

Las \glspl{cnn} parten de la idea de mirar la imagen por partes pequeñas en lugar de toda junta. La operación
que lo hace es la \textbf{convolución} (\emph{convolution}), y su herramienta es un filtro (\emph{kernel}): una
pequeña matriz de pesos, por ejemplo de $3 \times 3$, que codifica un patrón visual, como el paso de un tono
claro a uno oscuro. El filtro se apoya sobre una ventana de $3 \times 3$ píxeles, multiplica cada peso por el
píxel que tiene debajo y suma los nueve productos, es decir, calcula el producto punto entre sus pesos y la
ventana. Cuanto más se parece la ventana al patrón que codifica el filtro, mayor es el resultado. Por ejemplo,
un filtro con pesos positivos en su columna izquierda, nulos en la central y negativos en la derecha da un valor
grande sobre una ventana clara a la izquierda y oscura a la derecha, porque los píxeles claros, que tienen
valores mayores que los oscuros, se multiplican por los pesos positivos. Sobre una ventana de tono uniforme, en
cambio, las dos columnas se cancelan y el resultado es cero. Después, el filtro se desplaza un píxel y repite el cálculo, hasta recorrer toda la imagen. Los resultados, ordenados según la
posición en que se calcularon, forman un mapa de características (\emph{feature map}) que indica dónde aparece
el patrón. En una imagen dermatoscópica, ese filtro produce valores grandes en los tramos verticales del contorno
de la lesión, donde la piel pasa de clara a oscura, y valores cercanos a cero sobre la piel sana, de tono
uniforme.

% Figura: la operación de convolución.
% Se usa en sec:marco_cnn, después del párrafo que define la convolución.
% Las tres matrices van separadas y las etiquetas en dos líneas, para que no se toquen (tutor, p36).

\begin{figure}[H]
    \centering
    \begin{tikzpicture}[x=1mm, y=1mm]

        % Entrada 6x6 con la ventana de 3x3 resaltada.
        \fill[fignaranja!30] (5, 10) rectangle (20, 25);
        \draw[figtrazo!50] (0, 0) grid[step=5] (30, 30);
        \draw[figtrazo, thick] (0, 0) rectangle (30, 30);
        \draw[fignaranja!70!black, very thick] (5, 10) rectangle (20, 25);
        \node[font=\footnotesize, text=figtrazo, align=center, anchor=north] at (15, -2) {entrada\\$6\times 6$};

        % Filtro 3x3.
        \fill[fignaranja!30] (50, 7.5) rectangle (65, 22.5);
        \draw[fignaranja!70!black, shift={(50, 7.5)}] (0, 0) grid[step=5] (15, 15);
        \draw[fignaranja!70!black, very thick] (50, 7.5) rectangle (65, 22.5);
        \node[font=\footnotesize, text=figtrazo, align=center, anchor=north] at (57.5, -2) {filtro $3\times 3$\\(pesos compartidos)};
        \node[font=\large, text=figtrazo] at (40, 15) {$\ast$};

        % Salida 4x4 con la celda resultante.
        \fill[figverde!35] (95, 10) rectangle (100, 15);
        \draw[figtrazo!50] (90, 0) grid[step=5] (110, 20);
        \draw[figtrazo, thick] (90, 0) rectangle (110, 20);
        \draw[figverde!70!black, very thick] (95, 10) rectangle (100, 15);
        \node[font=\footnotesize, text=figtrazo, align=center, anchor=north] at (100, -2) {mapa de características\\$4\times 4$};
        \node[font=\large, text=figtrazo] at (77.5, 15) {$=$};

        % De la ventana a la celda.
        \draw[figflecha, fignaranja!70!black, bend left=25] (20, 25) to
            node[figetiqueta, above] {producto punto} (95, 15);

    \end{tikzpicture}
    \caption{Convolución: un mismo filtro recorre la entrada y produce un mapa de características.}
    \label{fig:marco_convolucion}
\end{figure}


La Figura~\ref{fig:marco_convolucion} muestra el recorrido sobre una entrada de $6 \times 6$: como el filtro
cabe en cuatro posiciones por fila y cuatro por columna, el mapa es de $4 \times 4$. La operación responde a los dos rasgos de las imágenes señalados antes. Como cada valor del mapa depende solo de
una ventana pequeña, la red supone que la información relevante es local, y como el mismo filtro, con los mismos
pesos, se aplica en toda la imagen, el contorno de la lesión se detecta en cualquier lugar en que aparezca. Los
pesos de los filtros no se fijan a mano, sino que se aprenden durante el entrenamiento, como los de cualquier
otra capa. Una capa convolucional aplica muchos filtros a la vez, cada uno con sus propios pesos, y apila sus
mapas como canales de la salida~\cite{lecun1998gradient,goodfellow2016deep}.

Una \gls{cnn} completa alterna capas convolucionales, funciones de activación y capas de reducción
(\emph{pooling}), que resumen cada ventana de un mapa en un solo valor, por ejemplo su máximo, y reducen así la
resolución. Cada reducción hace que una unidad de la capa siguiente dependa de una región más grande de la
imagen original, llamada su campo receptivo (\emph{receptive field}). De este modo, la red concreta la cadena
de rasgos simples a complejos descrita al comienzo de la Sección~\ref{sec:marco_fundamentos}: las primeras capas
responden a bordes y las últimas, a partes de objetos~\cite{goodfellow2016deep}. En una imagen dermatoscópica,
esa cadena iría del contorno y los cambios de color a estructuras que ocupan parte de la lesión, como una red de
líneas pigmentadas o un grupo de puntos~\cite{argenziano1998checklist}, y de ellas a la lesión completa. Al
final de la red, el
promedio de cada canal sobre todas las posiciones convierte los mapas en un vector de características, y la
cabeza de clasificación de la Sección~\ref{sec:marco_transferencia} lo convierte en los
logits~\cite{he2015residual}.

Agregar capas, sin embargo, no siempre mejora la red. A partir de cierta profundidad aparece la degradación
(\emph{degradation}): el error aumenta al sumar capas. Por ejemplo, una red convolucional de 56 capas obtiene
más error de entrenamiento que una de 20 sobre el mismo conjunto~\cite{he2015residual}. Como el error crece
sobre los propios datos de entrenamiento, el problema no es el sobreajuste de la
Sección~\ref{sec:marco_regularizacion}, sino la dificultad de optimizar la red. En principio, a la red de 56
capas le bastaría con que sus 36 capas adicionales calcularan la identidad, es decir, devolvieran su entrada
sin cambios, para igualar a la de 20. Sin embargo, al optimizador le resulta difícil aproximar la identidad con
varias capas no lineales~\cite{he2015residual}.

ResNet resolvió esta dificultad con las \textbf{conexiones residuales} (\emph{residual connections}). Una
conexión residual suma la entrada de un bloque de capas al resultado de esas capas, de modo que la entrada
llega a la salida del bloque también por un camino que no la transforma. Así, el bloque ya no tiene que producir su salida completa, sino solo la
corrección $F$ que hay que agregarle a la entrada,
\begin{equation}
    \mathbf{h}^{(l+1)} = \mathbf{h}^{(l)} + F\bigl(\mathbf{h}^{(l)}\bigr),
    \label{eq:marco_residual}
\end{equation}
donde $\mathbf{h}^{(l)}$ es la entrada del bloque $l$. Si lo mejor para un bloque es no modificar su entrada, le
basta con llevar los pesos de $F$ a cero, y su salida es igual a su entrada. De este modo, en el peor caso un
bloque nuevo no cambia nada, y en el mejor aporta una corrección útil. Con esta modificación se pudieron entrenar redes de más de cien capas~\cite{he2015residual}, y la misma conexión
reaparece en el transformer (Sección~\ref{sec:marco_atencion}).

Las dos redes convolucionales que evalúa el trabajo parten de esta idea y la desarrollan en direcciones
distintas. DenseNet conecta cada capa con todas las anteriores de su bloque: cada capa recibe, concatenados, los
mapas de características de las capas previas, de modo que reutiliza los rasgos ya calculados y alcanza una
exactitud comparable a la de una ResNet con menos parámetros~\cite{huang2018densely}. ConvNeXt, en cambio,
conserva la estructura de una ResNet y actualiza su diseño con decisiones tomadas de los transformers que se
presentan a continuación, como filtros de $7 \times 7$ y la normalización por capa~\cite{liu2022convnet}.
Aun con estas mejoras, una convolución solo relaciona píxeles cercanos. Para vincular dos regiones alejadas de
la imagen, por ejemplo, para comparar un borde de la lesión con el opuesto, la red necesita apilar capas hasta
que su campo receptivo cubra a ambos.

\subsection{Atención y arquitectura transformer}
\label{sec:marco_atencion}

La \textbf{atención} (\emph{attention}), en cambio, relaciona cualquier par de posiciones de la entrada en una
sola capa, sin importar la distancia entre ellas. Se propuso en 2017 para la traducción de textos, como núcleo
de una arquitectura llamada transformer~\cite{vaswani2017attention}, y hoy los modelos de lenguaje y buena
parte de los modelos de imagen se construyen sobre ella~\cite{dosovitskiy2021image}. Para aplicarla, la entrada
se divide en una secuencia de elementos llamados \emph{tokens}. En un texto, cada token es una palabra o
un fragmento de palabra, y en una imagen, un parche (\emph{patch}), es decir, un pequeño recuadro de píxeles.
Cada token se representa con un \emph{embedding}, un vector numérico que el modelo aprende de modo que
los tokens con significados parecidos tengan vectores parecidos.

La atención determina, para cada token, qué información tomar de los demás. Para ello, proyecta el
embedding de cada token en tres vectores, con matrices de pesos aprendidas: una consulta
(\emph{query}), que expresa qué información busca el token; una clave (\emph{key}), que expresa qué
información ofrece a los demás, y un valor (\emph{value}), que es la información que entrega. La compatibilidad
entre dos tokens se mide con el producto punto entre la consulta de uno y la clave del otro, y la
softmax de la Ecuación~\ref{eq:marco_softmax} convierte esas compatibilidades en pesos que suman uno. La salida
de cada token es el promedio de los valores de todos los tokens, ponderado por esos pesos. Si se
reúnen las consultas, las claves y los valores en las matrices $Q$, $K$ y $V$, la operación completa, llamada
atención escalada por producto punto, es
\begin{equation}
    \mathrm{Attention}(Q, K, V) = \mathrm{softmax}\!\left(\frac{QK^{\top}}{\sqrt{d_k}}\right) V,
    \label{eq:marco_atencion}
\end{equation}
donde $d_k$ es la dimensión de las claves. La división por $\sqrt{d_k}$ evita que los productos punto crezcan
con la dimensión y lleven la softmax a zonas donde su gradiente es casi nulo, lo que dificultaría el
entrenamiento~\cite{vaswani2017attention}. La Figura~\ref{fig:marco_atencion} muestra los pesos para una
secuencia de cuatro tokens: cada fila corresponde a una consulta, y sus pesos indican cuánto aporta el
valor de cada posición a la salida de ese token.

% Figura: la matriz de pesos de atención.
% Se usa en sec:marco_atencion, después de la Ecuación de la atención escalada por producto punto.
% Los pesos son ilustrativos: muestran la forma de la operación, no un modelo entrenado.

\begin{figure}[H]
    \centering
    \begin{tikzpicture}[x=1mm, y=1mm]

        % Consultas (filas) y claves (columnas).
        \foreach \i in {1, 2, 3, 4} {
            \node[fignodo, minimum height=6mm, minimum width=9mm, inner sep=1pt,
                  draw=figazul!80!black, fill=figazul!12] (q\i) at (0, 30 - \i*8) {\footnotesize $q_{\i}$};
            \node[fignodo, minimum height=6mm, minimum width=6.5mm, inner sep=1pt,
                  draw=fignaranja!70!black, fill=fignaranja!18] (k\i) at (8 + \i*8, 34) {\footnotesize $k_{\i}$};
        }

        % Pesos softmax(QK^T / sqrt(d_k)); cada fila suma uno.
        \foreach \row/\vals in {1/{60,20,10,10}, 2/{15,55,20,10}, 3/{10,25,50,15}, 4/{5,10,25,60}} {
            \foreach \v [count=\col] in \vals {
                \fill[figceleste!\v!white] (4 + \col*8, 26 - \row*8) rectangle (12 + \col*8, 34 - \row*8);
                \node[font=\scriptsize] at (8 + \col*8, 30 - \row*8) {0.\ifnum\v<10 0\fi\v};
            }
        }
        \draw[figtrazo, shift={(12, -6)}] (0, 0) grid[step=8] (32, 32);

        % Salida.
        \foreach \i in {1, 2, 3, 4} {
            \node[fignodo, minimum height=6mm, minimum width=9mm, inner sep=1pt,
                  draw=figverde!70!black, fill=figverde!10, densely dashed] (o\i) at (66, 30 - \i*8)
                  {\footnotesize $o_{\i}$};
            \draw[figflecha] (45, 30 - \i*8) -- (o\i.west);
        }

        \node[font=\footnotesize, text=figtrazo] at (66, 32) {$o_i = \sum_j a_{ij} v_j$};
        \node[font=\footnotesize, text=figtrazo, align=center] at (28, -12)
            {pesos $a_{ij}$ = $\mathrm{softmax}\!\left(q_i k_j^{\top} / \sqrt{d_k}\right)$};

    \end{tikzpicture}
    \caption{Pesos de atención entre las posiciones de una secuencia de cuatro tokens (valores ilustrativos).}
    \label{fig:marco_atencion}
\end{figure}


Un solo cálculo de atención reúne en un único conjunto de pesos todas las relaciones entre tokens. La
atención multicabezal (\emph{multi-head attention}) realiza varios cálculos en paralelo, cada uno con sus
propias proyecciones de menor dimensión, y concatena los resultados, de modo que cada cabeza puede capturar un
tipo distinto de relación entre posiciones~\cite{vaswani2017attention}.

La atención, tal como se definió, no tiene en cuenta el orden de los tokens: si se reordena la entrada,
cada token obtiene el mismo resultado que antes. Para que el modelo distinga las posiciones, se suma al
embedding de cada token un vector que depende de su lugar en la secuencia, llamado codificación de
posición (\emph{positional encoding}), que puede ser fijo o aprenderse durante el
entrenamiento~\cite{vaswani2017attention,dosovitskiy2021image}.

Un transformer apila bloques iguales. Cada bloque aplica primero una capa de atención multicabezal, que combina
información entre tokens, y después un perceptrón multicapa, que transforma cada token por
separado. Las dos partes llevan una conexión residual como la de la Ecuación~\ref{eq:marco_residual} y una
normalización por capa (\emph{layer normalization}), que lleva los valores de cada token a una escala
común~\cite{vaswani2017attention}. El transformer original combina dos pilas de bloques: un codificador
(\emph{encoder}), que procesa la entrada completa y produce una representación de cada token que
incorpora información de los demás, y un decodificador (\emph{decoder}), que genera la salida de a un
token por vez y se presenta en la Sección~\ref{sec:marco_llm}. El codificador no depende de qué
representen los tokens, de modo que también puede procesar una imagen si sus tokens son parches.

\subsection{Vision Transformer}
\label{sec:marco_vit}

Hasta la aparición del \gls{vit}, la atención se usaba en visión junto con las \glspl{cnn} o en reemplazo de
algunas de sus partes, y la convolución seguía siendo la base de los clasificadores. El \gls{vit} probó si un
transformer sin convoluciones, aplicado directamente a la imagen, podía clasificarla~\cite{dosovitskiy2021image}.
Para ello, aplica el codificador del transformer y convierte la imagen en una secuencia de tokens
(Figura~\ref{fig:marco_vit_patches}). Primero divide la imagen en una
rejilla de parches cuadrados, por ejemplo de $16 \times 16$ píxeles, de modo que una imagen de
$224 \times 224$ píxeles da 196 parches. En una imagen dermatoscópica, cada parche cubre un fragmento de la
lesión o de la piel que la rodea. Después aplana cada parche en un vector y lo proyecta, con una capa
lineal, a un embedding de dimensión fija, al que suma una codificación de posición aprendida que indica
su lugar en la rejilla. Al comienzo de la secuencia agrega un token de clase, que no corresponde a ningún
parche. Al atravesar los bloques del codificador, ese token reúne información de todos los parches, y la
cabeza de clasificación toma su representación final para producir los
logits~\cite{dosovitskiy2021image}.

% Figura 2.1. Del parche al token en el Vision Transformer.
% Se usa en sec:marco_vision, subsección "Vision Transformer y Swin Transformer", después del
% párrafo que define el ViT. Cada parche de la fila superior y su token comparten color, para que se
% lea la correspondencia; el token de clase va en gris porque no viene de ningún parche.

\begin{figure}[H]
    \centering
    \begin{tikzpicture}[x=1mm, y=1mm]

        % Imagen dividida en una rejilla de 4x4 parches; la fila superior va coloreada.
        \draw[figtrazo, thick, fill=figrelleno] (-26, -12) rectangle (0, 12);
        \foreach \c [count=\i from 0] in {figazul, fignaranja, figverde, figpurpura} {
            \fill[\c!35] (-26 + \i*6.5, 6) rectangle (-19.5 + \i*6.5, 12);
        }
        \foreach \x in {-19.5, -13, -6.5} { \draw[figtrazo] (\x, -12) -- (\x, 12); }
        \foreach \y in {-6, 0, 6} { \draw[figtrazo] (-26, \y) -- (0, \y); }
        \draw[figtrazo, thick] (-26, -12) rectangle (0, 12);
        \node[font=\small, text=figtrazo, align=center] at (-13, -18) {Imagen dividida\\en parches};

        % Operación de proyección.
        \node[figtarea, text width=28mm] (proj) at (22, 0)
            {Proyectar cada parche\\a un token};

        % Token de clase y tokens de los parches, con su embedding de posición.
        \node[fignodo, minimum width=7mm, minimum height=6mm, inner sep=1pt, fill=figrelleno]
            (t0) at (50, 6) {\footnotesize[CLS]};
        \node[font=\footnotesize, text=figtrazo] at (50, -3) {$+$};
        \node[font=\footnotesize, text=figtrazo] at (50, -9) {$p_{0}$};
        \foreach \c [count=\i from 1] in {figazul, fignaranja, figverde, figpurpura} {
            \node[fignodo, draw=\c!70!black, fill=\c!35, minimum width=6mm, minimum height=6mm,
                  inner sep=1pt] (t\i) at (50 + \i*9, 6) {\footnotesize $x_{\i}$};
            \node[font=\footnotesize, text=figtrazo] at (50 + \i*9, -3) {$+$};
            \node[font=\footnotesize, text=figtrazo] at (50 + \i*9, -9) {$p_{\i}$};
        }
        \node[font=\small, text=figtrazo] at (95, 6) {$\cdots$};
        \node[font=\small, text=figtrazo, align=center] at (72, -18)
            {Secuencia de tokens con\\embedding de posición por parche};

        % Flechas.
        \draw[figflecha] (0, 0) -- (proj.west);
        \draw[figflecha] (proj.east) -- (43, 0) |- (t0.west);

    \end{tikzpicture}
    \caption{Del parche al token en el Vision Transformer, con el token de clase ([CLS]) al inicio.}
    \label{fig:marco_vit_patches}
\end{figure}


Esta forma de procesar la imagen supone mucho menos sobre ella que una \gls{cnn}. Los supuestos sobre los datos
que una arquitectura incorpora de antemano, antes de ver ningún ejemplo, se llaman \textbf{sesgo inductivo}
(\emph{inductive bias}). En una \gls{cnn}, la localidad y la equivariancia a traslaciones, es decir, que un
patrón desplazado produzca la misma respuesta desplazada, están presentes en cada capa. En el \gls{vit}, en
cambio, solo el perceptrón de cada bloque es local, porque opera sobre cada token por separado, mientras
que la atención relaciona todos los parches desde la primera capa. Por lo tanto, las relaciones espaciales entre
parches tienen que aprenderse de los datos, y el \gls{vit} necesita más ejemplos que una \gls{cnn}: preentrenado
solo sobre \emph{ImageNet}, alcanza una exactitud algo menor que la de \glspl{cnn} de tamaño comparable, y
preentrenado sobre conjuntos de entre 14 y 300 millones de imágenes, las supera~\cite{dosovitskiy2021image}. En
consecuencia, el desempeño de un \gls{vit} depende en gran medida del preentrenamiento del que parte antes del
ajuste fino (Sección~\ref{sec:marco_transferencia}).

\subsection{Swin Transformer}
\label{sec:marco_swin}

El \gls{vit} mostró que la atención alcanza para clasificar imágenes, pero dejó dos problemas. El primero es el
costo: la atención global crece con el cuadrado del número de parches, porque relaciona cada token con todos
los demás, y eso limita la resolución de las imágenes que se pueden procesar, ya que duplicar el lado de una
imagen cuadruplica la cantidad de parches. El segundo es que procesa toda la imagen a una sola escala: sus
parches conservan el mismo tamaño en todos los bloques, sin la pirámide de mapas cada vez más pequeños de una
\gls{cnn}, que sirve cuando los objetos aparecen a tamaños distintos~\cite{liu2021swin}. El Swin Transformer
resuelve ambos problemas.

Para reducir el costo, el Swin Transformer limita la atención a ventanas locales de $M \times M$ parches. Si el
mapa tiene $h \times w$ parches y $D$ canales, el costo de la atención global y el de la atención por ventanas
son
\begin{equation}
    \Omega_{\text{global}} = 4hwD^2 + 2(hw)^2D, \qquad \Omega_{\text{ventanas}} = 4hwD^2 + 2M^2hwD,
    \label{eq:marco_swin_costo}
\end{equation}
de modo que, con $M$ fijo, el segundo crece de forma lineal con el número de parches~\cite{liu2021swin}.

% Figura: ventanas y ventanas desplazadas del Swin Transformer. Adaptada de liu2021swin (Figura 2).
% Se usa en sec:marco_vit, en el párrafo del Swin.

\begin{figure}[H]
    \centering
    \begin{tikzpicture}[x=1mm, y=1mm]

        % Capa l: cuatro ventanas de 4x4 parches.
        \foreach \c/\x/\y in {figazul/0/16, fignaranja/16/16, figverde/0/0, figpurpura/16/0} {
            \fill[\c!22] (\x, \y) rectangle (\x + 16, \y + 16);
        }
        \draw[figtrazo!40] (0, 0) grid[step=4] (32, 32);
        \draw[figtrazo, very thick] (0, 0) rectangle (32, 32);
        \draw[figtrazo, very thick] (16, 0) -- (16, 32) (0, 16) -- (32, 16);
        \node[font=\footnotesize, text=figtrazo] at (16, -5) {capa $l$: ventanas regulares};

        % Capa l+1: ventanas desplazadas dos parches.
        \begin{scope}[xshift=50mm]
            \foreach \c/\x/\y/\w/\h in {
                figazul/0/24/8/8, fignaranja/8/24/16/8, figverde/24/24/8/8,
                figpurpura/0/8/8/16, figceleste/8/8/16/16, figbermellon/24/8/8/16,
                figazul/0/0/8/8, fignaranja/8/0/16/8, figverde/24/0/8/8} {
                \fill[\c!22] (\x, \y) rectangle (\x + \w, \y + \h);
            }
            \draw[figtrazo!40] (0, 0) grid[step=4] (32, 32);
            \draw[figtrazo, very thick] (0, 0) rectangle (32, 32);
            \draw[figtrazo, very thick] (8, 0) -- (8, 32) (24, 0) -- (24, 32) (0, 8) -- (32, 8) (0, 24) -- (32, 24);
            \node[font=\footnotesize, text=figtrazo] at (16, -5) {capa $l+1$: ventanas desplazadas};
        \end{scope}

        \draw[figflecha] (36, 16) -- (46, 16);

    \end{tikzpicture}
    \caption{Ventanas y ventanas desplazadas del Swin Transformer. Adaptada de \textcite{liu2021swin}.}
    \label{fig:marco_swin}
\end{figure}


Las ventanas fijas, sin embargo, no intercambian información entre sí. Para que la información pase de una
ventana a otra, los bloques alternan una partición regular en ventanas con otra desplazada en medio tamaño de
ventana, cuyas ventanas abarcan parches de varias ventanas de la partición anterior
(Figura~\ref{fig:marco_swin}). Para resolver el segundo problema, entre una etapa y la siguiente, una capa de
fusión de parches (\emph{patch merging})
combina cada grupo de $2 \times 2$ parches vecinos en uno solo, de modo que la resolución se reduce a la mitad
y la cantidad de canales se duplica. El resultado es una pirámide de mapas de características como la de una
\gls{cnn}. Por último, la posición de los parches no se suma al embedding como en el \gls{vit}, sino
que entra como un término que depende de la posición relativa entre cada par de parches de la ventana y se suma
a su compatibilidad~\cite{liu2021swin}.

Las arquitecturas presentadas en esta sección se diferencian, en definitiva, por cuánto suponen sobre la
imagen. Las \glspl{cnn} incorporan la localidad y la jerarquía en su diseño; el \gls{vit} prescinde de esos
supuestos y compensa su ausencia con más datos de preentrenamiento, y el Swin Transformer recupera la localidad
y la jerarquía con una operación de atención. La diferencia entre familias, sin embargo, no reside solo en la
operación de base: ConvNeXt, una red puramente convolucional con un diseño actualizado, alcanza en
\emph{ImageNet} una exactitud igual o superior a la del Swin Transformer~\cite{liu2022convnet}. Por eso el
trabajo evalúa representantes de ambas familias.

Estos clasificadores cumplen, además, un papel de referencia. Una \textbf{línea base} (\emph{baseline}) es un
modelo establecido contra el que se compara una propuesta nueva, porque un resultado aislado no permite juzgar
el desempeño de un modelo hasta que se lo contrasta con lo que logra la referencia sobre los mismos datos. Las
\glspl{cnn} y los transformers visuales entrenados con las etiquetas de la tarea son la línea base de los
\glspl{vlm} (Sección~\ref{sec:marco_vlm}): son las arquitecturas consolidadas para clasificar imágenes, y un
modelo de otra familia justifica su costo si iguala su desempeño o si aporta algo que ellas no ofrecen. Sin
embargo, ninguno de estos clasificadores produce texto, ya que su salida es una probabilidad por clase. Para
redactar un reporte hace falta un modelo capaz de generar secuencias de palabras, y los modelos de lenguaje
actuales se construyen sobre el decodificador del transformer.

\section{Modelos de lenguaje y adaptación eficiente}
\label{sec:marco_llm}

Un modelo de lenguaje asigna una probabilidad a cada continuación posible de un texto: dado lo escrito
hasta el momento, predice el token siguiente. Con el transformer, estos modelos crecieron en tamaño y en
datos hasta volverse capaces de seguir instrucciones escritas.

\subsection{Grandes modelos de lenguaje}
\label{sec:marco_llm_grandes}

Un \gls{llm} se obtiene por preentrenamiento autosupervisado sobre corpus de miles de millones de
palabras. La tarea es predecir el token siguiente, y la supervisión sale del propio texto, sin etiquetas
manuales. El modelo aprende así la probabilidad de una secuencia como el producto de las probabilidades de
cada token dados los anteriores,
\begin{equation}
    p(x_1, \dots, x_T) = \prod_{t=1}^{T} p(x_t \mid x_1, \dots, x_{t-1}),
    \label{eq:marco_lm}
\end{equation}
y de esa fase sale una distribución de probabilidad sobre el vocabulario condicionada a cualquier
prefijo~\cite{brown2020language,touvron2023llama}.

Según qué parte del transformer conserven, se distinguen dos familias. Los modelos de solo codificador,
como BERT, de 2018, construyen una representación bidireccional del texto, en la que cada token ve a los
anteriores y a los posteriores, y sirven para clasificar y para comparar textos~\cite{devlin2019bert}. Los de
solo decodificador ven únicamente los tokens anteriores y sirven para generar texto. En 2020 se observó que la
pérdida de estos modelos baja de forma predecible al aumentar sus parámetros, sus datos y su
cómputo~\cite{kaplan2020scaling}, y GPT-3, con 175 mil millones de parámetros, mostró las capacidades que
aparecen con esa escala~\cite{brown2020language}. Desde 2023 hay además familias de pesos abiertos, que se
pueden descargar y ajustar, como LLaMA~\cite{touvron2023llama} y Mistral, cuyo modelo de 7 mil millones de
parámetros supera a uno de LLaMA~2 de 13 mil millones~\cite{jiang2023mistral}.

La familia Phi siguió el camino opuesto al aumento de escala: reducir el tamaño y mejorar los datos. Sus
primeros modelos se entrenaron con textos de calidad de libro de texto, en parte
sintéticos~\cite{gunasekar2023textbooks}, y Phi-3-mini, con 3800 millones de parámetros, rinde a la par de
modelos mucho mayores y cabe en un teléfono. Su variante Phi-3.5-vision, de 4200 millones de parámetros,
agrega entrada visual~\cite{abdin2024phi3}. Todos estos son modelos de solo decodificador.

El comportamiento del modelo se ajusta después del preentrenamiento. El ajuste fino supervisado lo entrena
sobre pares instrucción--respuesta, y la alineación con \gls{rlhf} lo entrena a preferir las
respuestas que un evaluador humano marca como mejores. Un modelo alineado responde a instrucciones, de modo
que se le puede pedir una tarea escribiéndola en lenguaje natural~\cite{ouyang2022training}.

Todo el texto que el modelo recibe antes de generar, la instrucción y cualquier información adicional,
forma su contexto y se llama \emph{prompt}. Como la probabilidad de la Ecuación~\ref{eq:marco_lm} está
condicionada en ese prefijo, cambiar el \emph{prompt} cambia lo que el modelo genera sin modificar sus
pesos. Los modelos grandes aprovechan además la información y los ejemplos incluidos en el \emph{prompt}, sin
ningún paso de entrenamiento, lo que se conoce como aprendizaje en contexto
(\emph{in-context learning})~\cite{brown2020language}.

Dado un \emph{prompt}, el texto se genera token a token: en cada paso el modelo devuelve una distribución
sobre el vocabulario, se elige un token y se lo agrega al prefijo. La regla con que se elige es la
estrategia de decodificación. La decodificación voraz (\emph{greedy decoding}) toma siempre el token más
probable, y en la generación abierta, como los otros métodos que buscan el texto más probable, tiende a
repetirse~\cite{holtzman2020curious}. Los
métodos de muestreo (\emph{sampling}) sortean el token siguiente según la distribución del modelo. La
temperatura $T$ divide los logits antes de la softmax,
\begin{equation}
    p_i = \frac{\exp(z_i / T)}{\sum_{j} \exp(z_j / T)},
    \label{eq:marco_temperatura}
\end{equation}
de modo que con $T<1$ la probabilidad se concentra en los tokens más probables y con $T>1$ se reparte; es
la misma operación que la Sección~\ref{sec:marco_eval_calib} usa para calibrar un clasificador. El muestreo de
núcleo (\emph{top-p}) sortea entre el conjunto más chico de tokens cuya probabilidad acumulada alcanza un
umbral $p$, de modo que la cantidad de candidatos crece cuando el modelo duda y se achica cuando está seguro.
Con eso evita tanto la repetición de la decodificación voraz como la incoherencia de muestrear de la
distribución completa~\cite{holtzman2020curious}. Como la estrategia de decodificación cambia el texto
generado, forma parte de la configuración de un generador y hay que declararla al comparar modelos.

\subsection{Ajuste eficiente de parámetros: LoRA}
\label{sec:marco_llm_lora}

Adaptar uno de estos modelos a un dominio con el ajuste fino completo de la
Sección~\ref{sec:marco_transferencia} tropieza con la memoria: con Adam, cada parámetro entrenado arrastra su
gradiente y dos estados adicionales, y para miles de millones de parámetros eso puede superar la memoria de una
GPU. Los métodos de \gls{peft} congelan el modelo preentrenado y entrenan solo un conjunto pequeño de
parámetros, y reducen así la memoria que el ajuste necesita~\cite{hu2021lora}.

\Gls{lora}, propuesta en 2021, es uno de ellos. Parte de que el cambio que el ajuste fino introduce en una
matriz de pesos tiene rango bajo, es decir, se puede escribir como el producto de dos matrices delgadas. Sobre
cada matriz congelada $W$ que se elige adaptar entrena entonces un incremento de bajo rango,
\begin{equation}
    W' = W + \frac{\alpha}{r}\,BA, \qquad B \in \mathbb{R}^{d\times r},\quad A \in \mathbb{R}^{r\times k},\quad
    r \ll \min(d, k),
    \label{eq:marco_lora}
\end{equation}
y solo actualiza $A$ y $B$ (Figura~\ref{fig:marco_lora}). $A$ se inicializa con valores aleatorios y $B$ en
cero, de modo que el incremento vale cero al empezar y el modelo parte exactamente del preentrenado. El rango
$r$ fija cuántos parámetros tiene el incremento y con ello su capacidad de adaptación, y el factor $\alpha$
escala su efecto, de modo que se puede cambiar $r$ sin volver a ajustar la tasa de
aprendizaje~\cite{hu2021lora}.

% Figura: reparametrización de LoRA. Adaptada de hu2021lora (Figura 1).
% Se usa en sec:marco_llm_lora, después de la Ecuación de LoRA.

\begin{figure}[H]
    \centering
    \begin{tikzpicture}[x=1mm, y=1mm]

        \node[figentrada, text width=10mm, minimum height=7mm] (x) at (0, 0) {$\mathbf{x}$};

        % Rama congelada.
        \node[fignodo, text width=36mm, minimum height=24mm, fill=figrelleno] (W) at (-24, 25)
            {pesos preentrenados\\$W \in \mathbb{R}^{d\times k}$\\[2pt]{\footnotesize congelados}};

        % Rama entrenable: A reduce a r, B vuelve a d.
        \draw[draw=fignaranja!70!black, fill=fignaranja!25, thick] (8, 16) -- (32, 16) -- (26, 22) -- (14, 22) -- cycle;
        \node[font=\footnotesize] at (20, 19) {$A \in \mathbb{R}^{r\times k}$};
        \draw[draw=fignaranja!70!black, fill=fignaranja!25, thick] (14, 28) -- (26, 28) -- (32, 34) -- (8, 34) -- cycle;
        \node[font=\footnotesize] at (20, 31) {$B \in \mathbb{R}^{d\times r}$};
        \node[font=\footnotesize, text=fignaranja!70!black, align=left, anchor=west] at (34, 25)
            {entrenables\\rango $r \ll \min(d,k)$};

        % Suma y salida.
        \node[circle, draw=figtrazo, inner sep=1pt, font=\small] (sum) at (0, 46) {$+$};
        \node[figsalida, text width=26mm, minimum height=7mm] (h) at (0, 58) {$\mathbf{h} = W\mathbf{x} + BA\mathbf{x}$};

        \draw[figflecha] (x.north) -- ++(0, 4) -| (W.south);
        \draw[figflecha] (x.north) -- ++(0, 4) -| (20, 16);
        \draw[figflecha] (20, 22) -- (20, 28);
        \draw[figflecha] (W.north) |- (sum.west);
        \draw[figflecha] (20, 34) |- (sum.east);
        \draw[figflecha] (sum.north) -- (h.south);

    \end{tikzpicture}
    \caption{Reparametrización de LoRA con el incremento de bajo rango $BA$. Adaptada de \textcite{hu2021lora}.}
    \label{fig:marco_lora}
\end{figure}


El método no fija qué matrices adaptar. El trabajo original lo aplicó solo a las de la atención
(Sección~\ref{sec:marco_atencion}), en la mayoría de los experimentos a las proyecciones $Q$ y $V$, y dejó las
del perceptrón de cada bloque como trabajo futuro, aunque señala que se puede aplicar a cualquier subconjunto de
matrices. Cuantas más matrices se adaptan y mayor es
$r$, más parámetros se entrenan. Aun así son una fracción mínima: en GPT-3 redujo los parámetros entrenables
unas 10\,000 veces y la memoria de la GPU unas tres veces~\cite{hu2021lora}. Al terminar, el incremento se suma
a $W$ y el modelo resultante se ejecuta igual que el original, sin latencia adicional.

Con un \gls{llm} adaptable en una GPU de consumo, falta conectarlo con la imagen. Eso exige llevar los rasgos
de la imagen al espacio de tokens del modelo de lenguaje, o bien alinear imagen y texto en un espacio común.

\section{Modelos visión--lenguaje}
\label{sec:marco_vlm}

Un \gls{vlm} procesa imagen y texto en un mismo modelo, y hay dos formas de construirlo, que aparecieron en
ese orden. La primera, desde 2021, entrena un codificador de imagen y uno de texto para que proyecten al mismo
espacio, de modo que una imagen y su descripción queden cerca. El resultado sirve para clasificar y para
buscar casos parecidos, pero no redacta. La segunda, desde 2023, conecta un codificador visual a la entrada de
un \gls{llm}, que entonces genera texto a partir de la imagen. En ambas, el codificador de imagen suele ser un
\gls{vit} como los de la Sección~\ref{sec:marco_vit}~\cite{radford2021learning,liu2023visual}.

\subsection{Alineación contrastiva: CLIP}
\label{sec:marco_vlm_clip}

\Gls{clip} entrena un codificador de imagen y uno de texto sobre pares imagen--texto, cada uno con una
proyección final a un espacio común de embeddings. La similitud entre el embedding de una imagen ($\mathbf{u}$)
y el de un texto ($\mathbf{v}$) es el coseno del ángulo que forman,
\begin{equation}
    \cos(\mathbf{u}, \mathbf{v}) = \frac{\mathbf{u}^{\top}\mathbf{v}}{\lVert\mathbf{u}\rVert\,\lVert\mathbf{v}\rVert},
    \label{eq:marco_cos}
\end{equation}
que vale 1 si los vectores apuntan en la misma dirección, 0 si son ortogonales y $-1$ si son opuestos. En cada
lote de $N$ pares, el modelo calcula la similitud entre cada imagen y cada texto y aprende a distinguir los $N$
pares correctos entre las $N^2$ combinaciones posibles, lo que acerca los embeddings de cada par y aleja los de
pares distintos (Figura~\ref{fig:marco_clip}). El modelo original se entrenó así sobre 400 millones de pares
recogidos de internet, sin ninguna etiqueta de clase~\cite{radford2021learning}.

% Figura: entrenamiento contrastivo de CLIP. Adaptada de radford2021learning (Figura 1, panel 1).
% Se usa en sec:marco_vlm_clip, después del párrafo que define el objetivo contrastivo.

\begin{figure}[H]
    \centering
    \begin{tikzpicture}[x=1mm, y=1mm]

        % Codificador de texto (arriba) y sus embeddings T_1..T_4.
        \node[figentrada, text width=22mm, minimum height=7mm] (txt) at (-26, 44) {textos del lote};
        \node[figtarea, text width=30mm, minimum height=7mm] (enct) at (22, 44) {codificador de texto};
        \draw[figflecha] (txt.east) -- (enct.west);
        \foreach \j in {1,...,4} {
            \node[draw=figverde!70!black, fill=figverde!10, minimum width=9mm, minimum height=6mm,
                font=\footnotesize] (T\j) at (\j*10, 33) {$T_{\j}$};
        }
        \draw[figflecha] (enct.south) -- ++(0, -3.5);

        % Codificador de imagen (izquierda) y sus embeddings I_1..I_4.
        \node[figentrada, text width=18mm, minimum height=7mm] (img) at (-40, 15) {imágenes del lote};
        \node[figtarea, text width=18mm, minimum height=7mm] (enci) at (-17, 15) {codificador de imagen};
        \draw[figflecha] (img.east) -- (enci.west);
        \foreach \i in {1,...,4} {
            \node[draw=figazul!80!black, fill=figazul!12, minimum width=9mm, minimum height=6mm,
                font=\footnotesize] (I\i) at (0, 33-\i*8) {$I_{\i}$};
        }
        \draw[figflecha] (enci.east) -- (I2.west |- enci.east);

        % Matriz de similitudes; la diagonal son los pares correctos.
        \foreach \i in {1,...,4} {
            \foreach \j in {1,...,4} {
                \ifnum\i=\j
                    \node[draw=fignaranja!70!black, fill=fignaranja!30, thick, minimum width=9mm,
                        minimum height=6mm, font=\scriptsize] at (\j*10, 33-\i*8)
                        {$I_{\i}{\cdot}T_{\j}$};
                \else
                    \node[draw=figtrazo, fill=figrelleno, minimum width=9mm, minimum height=6mm,
                        font=\scriptsize, text=figtrazo] at (\j*10, 33-\i*8) {$I_{\i}{\cdot}T_{\j}$};
                \fi
            }
        }
        \node[figetiqueta, align=left, anchor=west] at (47, 17)
            {diagonal: pares\\correctos, se acercan};
        \node[figetiqueta, align=left, anchor=west] at (47, 5)
            {resto: pares\\distintos, se alejan};

    \end{tikzpicture}
    \caption{Entrenamiento contrastivo de CLIP. Adaptada de \textcite{radford2021learning}.}
    \label{fig:marco_clip}
\end{figure}


Ese espacio compartido permite clasificar sin entrenar un clasificador. Cada clase se escribe como un
texto, por ejemplo «una imagen dermatoscópica de un melanoma», y la imagen se asigna a la clase cuyo
embedding de texto le queda más cerca. Esta clasificación \emph{zero-shot}, que no usa ningún ejemplo
etiquetado de la tarea, depende de cómo se redacten esos textos: cambiar la plantilla del \emph{prompt}
cambia la exactitud~\cite{radford2021learning}.

Con etiquetas disponibles, el codificador de imagen se usa como extractor de características congelado, de dos
maneras. Un clasificador lineal entrenado sobre los embeddings, el \emph{linear probe}, mide cuánta información
sobre las clases contiene la representación~\cite{radford2021learning}. El método de los \gls{knn}, en cambio,
no entrena ningún parámetro. Con una base de $m$ casos etiquetados, para una imagen consulta $x$ busca los $k$
casos cuyo embedding $f(x_j)$ tiene mayor similitud coseno con el suyo,
\begin{equation}
    \mathcal{N}_k(x) = \operatorname*{arg\,max}_{S \subseteq \{1,\dots,m\},\ |S| = k}\ \sum_{j \in S}
    \cos\bigl(f(x),\, f(x_j)\bigr),
    \label{eq:marco_vecinos}
\end{equation}
y estima la probabilidad de cada clase $c$ como la proporción de vecinos que la tienen,
\begin{equation}
    \hat{p}_{\text{kNN}}(c \mid x) = \frac{1}{k}\sum_{j \in \mathcal{N}_k(x)} \mathbf{1}[\,y_j = c\,],
    \label{eq:marco_knn}
\end{equation}
donde $y_j$ es la etiqueta del caso $j$. La predicción es la clase de mayor proporción, el voto mayoritario,
que extiende a $k$ vecinos la regla del vecino más cercano~\cite{cover1967nearest}. Así, el espacio contrastivo sirve también de índice de casos parecidos,
que la Sección~\ref{sec:marco_rep_recuperacion} usa como evidencia para el texto.

Un modelo contrastivo solo relaciona lo que vio en sus pares de entrenamiento, y los pares de internet
casi no contienen imágenes dermatoscópicas. DermLIP, de 2025, aplica el mismo esquema a alrededor de un millón
de pares imagen--texto dermatológicos, y en clasificación \emph{zero-shot} de enfermedades de la piel supera a
los modelos biomédicos de propósito general~\cite{yan2025derm1miccv}. Un modelo contrastivo, sin embargo, no
produce texto nuevo, y para redactar hace falta conectar el codificador visual con un modelo de lenguaje
generativo.

\subsection{Modelos multimodales generativos}
\label{sec:marco_vlm_generativos}

LLaVA, de 2023, es un esquema representativo de los \glspl{vlm} generativos abiertos. Un codificador
visual del tipo de \gls{clip} extrae los rasgos de la imagen, y un proyector, en LLaVA una capa lineal, los
lleva a la dimensión de los embeddings de palabra del \gls{llm}. Los tokens visuales resultantes se concatenan
con la instrucción en la secuencia de entrada de un \gls{llm} de solo decodificador, que genera la
respuesta token a token (Figura~\ref{fig:marco_llava}). La imagen llega como una parte más de la secuencia de
entrada, sin la estructura codificador--decodificador del transformer original~\cite{liu2023visual}.

LLaVA se entrena en dos etapas. En la primera, con el codificador visual y el \gls{llm} congelados, solo se
ajusta el proyector sobre pares imagen--leyenda, para que los tokens visuales caigan en una región del
espacio que el \gls{llm} interprete. En la segunda, el \emph{visual instruction tuning}, se ajustan el
proyector y el \gls{llm} sobre conversaciones en las que se pide describir, razonar o responder preguntas
sobre una imagen~\cite{liu2023visual}. Como el \gls{llm} concentra la mayoría de los parámetros, en esa
segunda etapa se le puede aplicar \gls{lora} (Sección~\ref{sec:marco_llm_lora}) en lugar del ajuste completo,
mientras el codificador visual sigue congelado y el proyector se entrena entero.

% Figura: arquitectura de un VLM generativo tipo LLaVA. Adaptada de liu2023visual (Figura 1).
% Se usa en sec:marco_vlm_generativos, después del párrafo que presenta el esquema.
% [H]: como flotante saltaba a la página siguiente y cortaba el párrafo que venía después.

\begin{figure}[H]
    \centering
    \begin{tikzpicture}[x=1mm, y=1mm]

        % Rama visual.
        \node[figentrada, text width=18mm] (img) at (0, 0) {imagen};
        \node[figtarea, text width=22mm] (enc) at (0, 16) {codificador visual\\{\footnotesize (tipo CLIP)}};
        \node[figtarea, text width=22mm] (proj) at (0, 32) {proyector\\{\footnotesize (capa lineal)}};
        \draw[figflecha] (img) -- (enc);
        \draw[figflecha] (enc) -- (proj);

        % Rama de texto.
        \node[figentrada, text width=22mm] (instr) at (40, 0) {instrucción};
        \node[fignodo, text width=22mm, fill=figrelleno] (tok) at (40, 16) {tokenizador y\\\emph{embeddings}};
        \draw[figflecha] (instr) -- (tok);

        % Secuencia de entrada al LLM: cada grupo de tokens queda encima de la rama que lo produce, y un
        % marco agrupa sus casillas para que la flecha llegue al grupo y no a una casilla suelta.
        \foreach \k in {0,...,3} {
            \node[draw=figazul!80!black, fill=figazul!12, minimum width=4.5mm, minimum height=4.5mm]
                (v\k) at (-8.25+\k*5.5, 47) {};
            \node[draw=figtrazo, fill=figrelleno, minimum width=4.5mm, minimum height=4.5mm]
                (t\k) at (31.75+\k*5.5, 47) {};
        }
        \node[draw=figtrazo, densely dotted, rounded corners=1pt, inner sep=1.2mm, fit=(v0) (v3)] (vis) {};
        \node[draw=figtrazo, densely dotted, rounded corners=1pt, inner sep=1.2mm, fit=(t0) (t3)] (txt) {};
        \node[figetiqueta, anchor=east, align=right] at (vis.west) {tokens\\visuales};
        \node[figetiqueta, anchor=west, align=left] at (txt.east) {tokens\\de texto};
        \draw[figflecha] (proj.north) -- (vis.south);
        \draw[figflecha] (tok.north) -- (txt.south);

        % LLM y salida: los dos grupos entran como una sola secuencia.
        \node[figtarea, minimum width=66mm] (llm) at (20, 64) {LLM de solo decodificador};
        \draw[figflecha] (vis.north) -- (vis.north |- llm.south);
        \draw[figflecha] (txt.north) -- (txt.north |- llm.south);
        \node[figsalida, text width=30mm] (out) at (20, 80) {texto generado};
        \draw[figflecha] (llm) -- (out);

        \node[figetiqueta, text=fignaranja!70!black, align=left, anchor=west] at (62, 32)
            {etapa 1: se entrena\\solo el proyector};
        \node[figetiqueta, text=fignaranja!70!black, align=left, anchor=west] at (62, 64)
            {etapa 2: proyector\\y LLM};

    \end{tikzpicture}
    \caption{Arquitectura de un \gls{vlm} generativo del tipo de LLaVA. Adaptada de \textcite{liu2023visual}.}
    \label{fig:marco_llava}
\end{figure}


La misma receta se llevó después al dominio médico cambiando los datos. LLaVA-Med, de 2023, la repite con
figuras y leyendas de artículos biomédicos y se entrena en menos de quince horas con ocho
GPU~\cite{li2023llavamed}; su versión 1.5 usa Mistral como modelo de lenguaje~\cite{zeng2025mmskin}. SkinVL, de 2025, la repite con
conjuntos dermatológicos, entre ellos uno construido a partir de libros de texto de la
especialidad~\cite{zeng2025mmskin}. Phi-3.5-vision combina también un codificador de imagen con un modelo de
lenguaje de solo decodificador, en un tamaño menor~\cite{abdin2024phi3}.

Un \gls{vlm} generativo, en cualquiera de estas variantes, devuelve texto y no una distribución de probabilidad
sobre las clases. Para clasificar con él se puede buscar la clase en la respuesta generada o, sin generar,
calcular con la Ecuación~\ref{eq:marco_lm} la probabilidad que el modelo asigna al nombre de cada clase como
continuación del \emph{prompt}, normalizada por su cantidad de tokens para no favorecer los nombres cortos, y
elegir la mayor. Esos puntajes ordenan las clases, pero nada garantiza que estén
calibrados, y una respuesta generada no dice cuánto confiar en ella. Antes de pedirle un reporte hay que saber
cuándo una predicción merece afirmarse.

\section{Incertidumbre y predicción selectiva}
\label{sec:marco_incertidumbre}

Un clasificador devuelve siempre una categoría y una probabilidad por clase, aun ante una imagen que no sabe
distinguir. En medicina una respuesta equivocada presentada con la misma seguridad que una correcta puede
llevar a una conducta errónea, y reconocer cuándo no conviene responder es un problema con teoría propia.

\subsection{Incertidumbre en la clasificación diagnóstica}
\label{sec:marco_incert_fuentes}

La incertidumbre de una predicción tiene dos fuentes. La incertidumbre aleatoria proviene de los datos: una
imagen desenfocada, o una lesión cuyo aspecto es compatible con varias categorías, deja la respuesta
indeterminada aunque el modelo sea bueno, y reunir más datos no la reduce. La incertidumbre epistémica
proviene del modelo, que no vio suficientes ejemplos de cierto tipo de lesión, y disminuye con más
datos~\cite{kendall2017uncertainties}.

Una tercera fuente es el cambio de distribución (Sección~\ref{sec:marco_generalizacion}): cuando la imagen
proviene de un equipo, centro o población que el entrenamiento no cubrió, lo aprendido puede no aplicarse.
Como la softmax de la Ecuación~\ref{eq:marco_softmax} hace que las probabilidades sumen uno, no existe una
salida que exprese «este caso no lo conozco», y la probabilidad no refleja necesariamente la incertidumbre.

En dermatología los errores no cuestan lo mismo: no detectar una lesión maligna tiene consecuencias
distintas de derivar una benigna. Además, la decisión final la toma una persona. En un experimento con
dermatólogos y médicos de atención primaria, el apoyo de un sistema de aprendizaje profundo mejoró la
exactitud de los dos grupos~\cite{groh2024learningaided}, de modo que su salida es un insumo de la decisión
clínica y no la decisión misma. Un sistema que reconoce su límite y deriva puede ser preferible a uno que
siempre responde, y formalizar la abstención permite medir cuánto se gana.

\subsection{Predicción selectiva y abstención}
\label{sec:marco_pred_selectiva}

Un clasificador con opción de rechazo se describe con un par $(f, g)$. La función $f$ asigna una categoría a
cada imagen $x$, y la función $g$ decide si la predicción se acepta ($g(x) = 1$) o si el modelo se abstiene
($g(x) = 0$). La cobertura es la fracción de imágenes que el modelo acepta, y el riesgo selectivo es la tasa
de error medida solo sobre las aceptadas,
\begin{equation}
    \mathrm{cobertura}(g) = \mathbb{E}\bigl[g(x)\bigr], \qquad
    \mathrm{riesgo}(f, g) = \frac{\mathbb{E}\bigl[\mathbf{1}[f(x) \neq y]\; g(x)\bigr]}{\mathbb{E}\bigl[g(x)\bigr]},
    \label{eq:marco_riesgo_selectivo}
\end{equation}
donde $y$ es la categoría real. Sobre $n$ imágenes, la cobertura es la proporción de imágenes con
$g(x_i) = 1$ y el riesgo, la proporción de errores entre ellas~\cite{elyaniv2010foundations,geifman2017selective}.

Una forma habitual de construir $g$ es partir de una señal de confianza $\kappa(x)$, un número que crece con la
seguridad del modelo, y de un umbral $\tau$,
\begin{equation}
    g(x) = \mathbf{1}\bigl[\kappa(x) \ge \tau\bigr].
    \label{eq:marco_umbral_confianza}
\end{equation}
Subir $\tau$ reduce la cobertura, porque se aceptan menos imágenes, y reduce también el riesgo si la señal es
informativa, porque las imágenes descartadas son las más propensas al error. El compromiso se resume en la
curva riesgo--cobertura, que dibuja el riesgo selectivo contra la cobertura al recorrer todos los umbrales.
Una señal es mejor cuanto más rápido baja el riesgo al reducir la cobertura, y una que no distingue los
aciertos de los errores deja el riesgo en el nivel del error global. Todo depende entonces de qué señal se use.

\subsection{Señales de confianza y calibración}
\label{sec:marco_senales_confianza}

La señal más simple es la máxima probabilidad de la softmax, $\kappa(x) = \max_c p_c$, que sirve de línea
base para detectar predicciones erróneas~\cite{hendrycks2017baseline}. Para que un umbral sobre ella tenga una
lectura probabilística, el modelo tiene que estar calibrado, es decir, sus probabilidades tienen que coincidir con sus frecuencias de acierto: entre las
predicciones a las que asigna confianza $p$, acierta una fracción $p$. En la notación de
\textcite{guo2017calibration}, con $\hat{Y}$ la clase predicha y $\hat{P}$ su probabilidad,
$P(\hat{Y} = Y \mid \hat{P} = p) = p$. Las redes profundas actuales tienden a ser sobreconfiadas, y las más
profundas y anchas suelen estar peor calibradas que las pequeñas aunque clasifiquen
mejor~\cite{guo2017calibration}. Con una confianza calibrada, un umbral de 0.9 acepta predicciones con al
menos 0.9 de probabilidad de ser correctas; sin calibración el umbral no tiene esa lectura. Las métricas de
calibración y sus correcciones se presentan en la Sección~\ref{sec:marco_eval_calib}.

Hay señales que no dependen de la probabilidad de un único modelo. El desacuerdo entre modelos aproxima la
incertidumbre del sistema, un principio que explotan los ensambles. Un ensamble promedia las probabilidades de
varios modelos entrenados por separado, lo que se llama voto suave, y suele ser más exacto y estar mejor
calibrado que cada uno de ellos~\cite{lakshminarayanan2017simple}.

\phantomsection\label{sec:marco_concordancia}%
Dos predictores con sesgos distintos, por ejemplo un clasificador supervisado y uno basado en vecinos
(Sección~\ref{sec:marco_vlm_clip}), pueden coincidir o no sobre la misma imagen. La \textbf{concordancia} entre
sus predicciones $\hat{y}_1(x)$ y $\hat{y}_2(x)$ es
\begin{equation}
    a(x) = \mathbf{1}\bigl[\hat{y}_1(x) = \hat{y}_2(x)\bigr],
    \label{eq:marco_concordancia}
\end{equation}
y sirve como señal solo si los casos en que los predictores coinciden se equivocan menos que aquellos en que
no, $P(\text{error} \mid a = 1) < P(\text{error} \mid a = 0)$, condición que se verifica con datos. Como usa
las decisiones y no las probabilidades, no depende de que los modelos estén calibrados, pero no informa
cuando los dos comparten un error.

El \textbf{consenso entre los vecinos} recuperados mide qué tan unánime es el voto de la
Ecuación~\ref{eq:marco_knn}, $\kappa(x) = \max_c \hat{p}_{\text{kNN}}(c \mid x)$. Vale 1 cuando los $k$
vecinos coinciden y depende de la representación y de la cobertura de la base: si la base no contiene casos
parecidos, los $k$ vecinos existen igual y no son informativos.

La calibración no crea información que la señal no contiene: un modelo que asigna 0.7 a todas sus
predicciones y acierta el 70\,\% de ellas está perfectamente calibrado, y su confianza no separa los casos
fáciles de los difíciles. Por eso una señal se evalúa con la curva riesgo--cobertura y no solo con una métrica
de calibración. Con una señal disponible, la decisión además deja de ser binaria.

\subsection{Tres salidas: afirmar, ofrecer diferenciales o derivar}
\label{sec:marco_tres_salidas}

\begingroup En la práctica clínica derivar no es la única alternativa a afirmar: un clínico que duda entre dos
o tres diagnósticos los enuncia como diagnóstico diferencial. Un sistema puede entonces tener tres salidas
(Tabla~\ref{tab:marco_salidas}): afirmar un diagnóstico, ofrecer diferenciales, es decir un conjunto de
candidatos, o derivar el caso cuando no se alcanza la confianza necesaria. Con dos umbrales
$\tau_1 > \tau_2$ sobre la señal, la política es\displaywidowpenalty=10000
\begin{equation}
    \pi(x) =
    \begin{cases}
        \text{afirmar} & \text{si } \kappa(x) \ge \tau_1,\\
        \text{diferencial} & \text{si } \tau_2 \le \kappa(x) < \tau_1,\\
        \text{derivar} & \text{si } \kappa(x) < \tau_2.
    \end{cases}
    \label{eq:marco_politica}
\end{equation}\endgroup

\begin{table}[H]
    \centering
    \caption{Salidas de un sistema con abstención: qué entrega cada una, cuándo corresponde y qué arriesga.}
    \label{tab:marco_salidas}
    \small
    \begin{tabularx}{\textwidth}{@{}>{\raggedright\arraybackslash}p{2.9cm}
        >{\raggedright\arraybackslash\hspace{0pt}}X >{\raggedright\arraybackslash\hspace{0pt}}X
        >{\raggedright\arraybackslash\hspace{0pt}}X@{}}
        \toprule
        \textbf{Salida} & \textbf{Qué entrega} & \textbf{Cuándo corresponde} & \textbf{Qué arriesga} \\
        \midrule
        Afirmar & Una categoría diagnóstica & Confianza alta
            & Un error que el lector no detecta por la forma del texto \\
        \filasep
        Ofrecer diferenciales & Un conjunto de categorías candidatas & Confianza intermedia
            & Informa menos, porque deja la elección entre candidatos al lector \\
        \filasep
        Derivar & Ninguna categoría; el caso pasa a un especialista & Confianza baja
            & No produce diagnóstico y traslada el caso a una consulta \\
        \bottomrule
    \end{tabularx}
\end{table}

La salida diferencial es una predicción por conjuntos, $S(x) \subseteq \{1, \dots, C\}$. El conjunto puede
formarse con las clases más probables de un modelo, hasta que su probabilidad acumulada alcance un umbral, o con
las predicciones de los distintos predictores cuando no coinciden. Cuando $|S(x)| = 1$, la salida coincide con
afirmar. Cuanto más caro es un error frente a derivar, más alto conviene fijar el umbral de afirmación, y esa
elección se hace en validación, nunca sobre el test.

Esta política decide cuánto afirma el sistema, pero no redacta nada. Cuando la salida tiene que ser un reporte
clínico, la decisión pasa a condicionar un texto, y un texto tiene problemas propios.

\section{Generación de reportes}
\label{sec:marco_reportes}

\subsection{El reporte clínico dermatológico}
\label{sec:marco_rep_tarea}

Un reporte clínico dermatológico es el registro escrito del examen de una lesión. Lo redacta el médico que la
examina, queda en la historia clínica y comunica lo observado a otros profesionales, por ejemplo al
especialista al que se deriva el caso. Su contenido se organiza en tres partes
(Tabla~\ref{tab:marco_partes_reporte}): la descripción de los hallazgos visibles, la impresión diagnóstica que
los resume y, cuando corresponde, una recomendación de conducta. La descripción usa un vocabulario
estandarizado por consenso, con términos para el color, el patrón, la distribución y las estructuras de la
lesión, y algoritmos como la lista de 7 puntos convierten la presencia de ciertos hallazgos en peso
diagnóstico~\cite{kittler2016standardization,argenziano1998checklist}.

\begin{table}[H]
    \centering
    \caption{Partes de un reporte dermatológico y lo que se le exige a cada una.}
    \label{tab:marco_partes_reporte}
    \small
    \begin{tabularx}{\textwidth}{@{}>{\raggedright\arraybackslash}p{3.0cm}
        >{\raggedright\arraybackslash\hspace{0pt}}X >{\raggedright\arraybackslash\hspace{0pt}}X@{}}
        \toprule
        \textbf{Parte} & \textbf{Contenido} & \textbf{Qué se le exige} \\
        \midrule
        Descripción & Color, patrón, bordes, estructuras y localización de la lesión
            & Que describa lo que la imagen muestra, con el vocabulario de la especialidad \\
        \filasep
        Impresión diagnóstica & La categoría diagnóstica, o los diferenciales si hay duda
            & Que sea la correcta y que su grado de certeza corresponda a la evidencia \\
        \filasep
        Recomendación & Seguimiento, biopsia o derivación
            & Que sea coherente con la impresión diagnóstica \\
        \bottomrule
    \end{tabularx}
\end{table}

Generar un reporte de forma automática tiene una historia corta. Su antecedente es la descripción automática
de imágenes (\emph{image captioning}), que produce una oración sobre una imagen natural. El reporte médico es
más largo, sigue una estructura y usa un vocabulario técnico. Su generación con redes neuronales se
desarrolló primero en radiología, desde 2018, con modelos que predicen etiquetas de hallazgos y redactan el reporte oración por
oración, y con arquitecturas que recuerdan patrones de redacción entre reportes
parecidos~\cite{jing2018automatic,chen2020generating}. Con los \glspl{vlm} generativos, la tarea pasó a
resolverse con un único modelo que recibe la imagen y una instrucción~\cite{hartsock2024vision}. En
todos los casos el entrenamiento requiere pares de imagen y texto clínico, y en dermatología esos pares son
escasos y recientes: su texto fue redactado por dermatólogos, tomado de libros de la especialidad o generado
por un modelo y revisado por especialistas~\cite{shen2025skincare,zeng2025mmskin,yan2025derm1miccv}, y ese origen limita cuánto se puede
confiar en él como referencia.

\subsection{Tres ejes de la generación de reportes}
\label{sec:marco_rep_ejes}

Generar un reporte resuelve tres problemas a la vez, uno por cada parte importante del texto.

\looseness=-1 El primero es describir lo que se ve: el texto tiene que reflejar los rasgos visuales de la
lesión con el vocabulario de la especialidad. La \textbf{calidad descriptiva} es la correspondencia entre lo
que el texto dice y lo que la imagen muestra, y el vocabulario cerrado permite medirla
(Sección~\ref{sec:marco_eval_texto}).

El segundo es incorporar información diagnóstica, porque el texto tiene que nombrar una categoría. La
\textbf{fidelidad diagnóstica} es que el diagnóstico que nombra el texto sea el correcto, y es independiente de
la calidad descriptiva: un texto puede describir fielmente la imagen y cerrar en un diagnóstico equivocado, y
otro puede nombrar el diagnóstico correcto acompañado de hallazgos inventados. Un reporte es correcto solo si
cumple las dos condiciones (Figura~\ref{fig:intro_dos_tareas}).

% Las dos condiciones de un reporte clínico correcto (clasificar y describir una misma imagen).
% Se trajo desde el Capítulo 1 a sec:marco_rep_ejes (2.6.2). El nombre del archivo y la etiqueta
% (fig:intro_dos_tareas) se conservan.

\begin{figure}[H]
    \centering
    \begin{tikzpicture}[x=1mm, y=1mm]

        \node[figentrada, text width=27mm] (img) at (12, 0) {Imagen\\dermatoscópica};

        \node[figtarea, text width=36mm] (t1) at (58, 14)
            {\textbf{Clasificar}\\asignar una categoría};
        \node[figtarea, text width=36mm] (t2) at (58, -14)
            {\textbf{Describir}\\redactar el texto\\de la lesión};

        \node[figsalida, text width=34mm] (cat) at (114, 14)
            {Categoría diagnóstica\\[2pt]{\footnotesize se evalúa con la\\fidelidad diagnóstica}};
        \node[figsalida, text width=34mm] (txt) at (114, -14)
            {Descripción de la lesión\\[2pt]{\footnotesize se evalúa con la\\calidad descriptiva}};

        \draw[figflecha] (img.east) -- ++(6, 0) |- (t1.west);
        \draw[figflecha] (img.east) -- ++(6, 0) |- (t2.west);
        \draw[figflecha] (t1.east) -- (cat.west);
        \draw[figflecha] (t2.east) -- (txt.west);

        \begin{scope}[on background layer]
            \node[figgrupo={las dos condiciones de un reporte correcto},
                  fit=(t1) (t2)] {};
        \end{scope}

    \end{tikzpicture}
    \caption{Los dos requisitos de un reporte clínico correcto: la categoría diagnóstica que nombra se evalúa por su fidelidad diagnóstica, y la descripción de la lesión, por su calidad descriptiva.}
    \label{fig:intro_dos_tareas}
\end{figure}


El tercero es tratar las predicciones inciertas. Cuando el diagnóstico no está claro, afirmar uno arriesga un
error que el lector no puede detectar, y la política de la Sección~\ref{sec:marco_tres_salidas} decide si el
texto afirma, enumera diferenciales o no se redacta. Los tres problemas dependen de cómo se produce el texto.

\subsection{Cómo se genera un reporte}
\label{sec:marco_rep_generacion}

Hay dos formas de generar un reporte con un \gls{vlm} (Figura~\ref{fig:marco_generacion_reportes}). En la
generación \textbf{de extremo a extremo}, el modelo recibe la imagen $x$ y una instrucción, y produce el texto
completo, $y \sim p_\theta(\cdot \mid x)$: describe la lesión e infiere el diagnóstico a partir de lo que ve.
Resuelve los tres problemas con un solo modelo, que tiene que ser a la vez un buen clasificador y un buen
redactor.

% Figura: dos formas de generar un reporte con un VLM, de extremo a extremo o condicionado en evidencia
% externa. Se usa en sec:marco_rep_generacion.

\begin{figure}[H]
    \centering
    \begin{tikzpicture}[x=1mm, y=1mm]

        % (a) De extremo a extremo.
        \node[font=\small\bfseries, anchor=west] at (-6, 30) {(a) De extremo a extremo};
        \node[figentrada, text width=22mm] (a_img) at (8, 18) {Imagen e\\instrucción};
        \node[figtarea, text width=26mm] (a_vlm) at (62, 18) {\gls{vlm}};
        \node[figsalida, text width=26mm] (a_rep) at (122, 18) {Reporte};
        \draw[figflecha] (a_img) -- (a_vlm);
        \draw[figflecha] (a_vlm) -- (a_rep);

        % (b) Condicionado en evidencia externa. Arriba, el componente que aporta la evidencia y la regla de
        % decisión; abajo, el mismo recorrido de (a) con la evidencia entrando al VLM desde arriba.
        \node[font=\small\bfseries, anchor=west] at (-6, 4) {(b) Condicionado en evidencia externa};
        \node[figtarea, text width=28mm] (b_ev) at (38, -8) {Clasificador o\\recuperación};
        \node[figtarea, text width=34mm] (b_pol) at (84, -8)
            {Regla de decisión\\[1pt]{\footnotesize afirmar, diferencial\\o derivar}};
        \node[figentrada, text width=22mm] (b_img) at (8, -34) {Imagen};
        \node[figtarea, text width=26mm] (b_vlm) at (84, -34) {\gls{vlm}};
        \node[figsalida, text width=26mm] (b_rep) at (122, -34) {Reporte};
        \draw[figflecha] (b_img.north) |- (b_ev.west);
        \draw[figflecha] (b_ev) -- (b_pol);
        \draw[figflecha] (b_pol) -- node[figetiqueta, right=1mm] {evidencia $c(x)$} (b_vlm);
        \draw[figflecha] (b_img) -- (b_vlm);
        \draw[figflecha] (b_vlm) -- (b_rep);

    \end{tikzpicture}
    \caption{Dos formas de generar un reporte: el \gls{vlm} solo o con la evidencia de
    otro componente.}
    \label{fig:marco_generacion_reportes}
\end{figure}


En la generación \textbf{condicionada en evidencia externa}, otro componente aporta información que se agrega al
\emph{prompt} del generador,
\begin{equation}
    y \sim p_\theta\bigl(\,\cdot \mid x,\ c(x)\bigr),
    \label{eq:marco_rep_condicionada}
\end{equation}
donde $c(x)$ es esa evidencia serializada como texto. Puede ser el diagnóstico que predice un clasificador, los
diagnósticos de casos parecidos recuperados de una base (Sección~\ref{sec:marco_rep_recuperacion}) o los
candidatos que deja la política de la Ecuación~\ref{eq:marco_politica}. El generador conserva la imagen y
redacta la descripción, pero el diagnóstico le llega como dato. Anclar así el texto en evidencia externa se
conoce como \emph{grounding}, y es aprendizaje en contexto (Sección~\ref{sec:marco_llm_grandes}): el
generador no cambia sus pesos~\cite{brown2020language}.

La segunda forma separa los problemas: la fidelidad diagnóstica pasa a depender de la evidencia, y la calidad
descriptiva sigue a cargo del generador. Tiene dos costos. El reporte hereda los errores de la evidencia,
porque un generador que recibe un diagnóstico equivocado lo afirma con la misma fluidez que uno correcto, y el
sistema usa más de un modelo. La instrucción también forma parte de la configuración: puede pedir que el
modelo razone antes de concluir o que entregue solo el reporte, y con eso cambia lo que el texto afirma. En las
dos formas el texto puede fallar de maneras que una métrica de clasificación no ve.

\subsection{Fallas del texto generado}
\label{sec:marco_rep_fallas}

Una alucinación es un texto fluido cuyo contenido no está sustentado por la entrada, que en este caso es la
imagen. Como falla, el fenómeno se describió primero en la generación de lenguaje
natural~\cite{ji2023survey} y en la descripción de imágenes, donde los modelos nombran objetos que la imagen no contiene en parte por el peso de
los patrones del lenguaje~\cite{rohrbach2018object}, y se documentó después en los
\glspl{vlm}~\cite{bai2024hallucination}. En los \glspl{vlm} médicos, un estudio sistemático encontró que los
modelos ajustados al dominio son más propensos a alucinar que los de propósito general, aun cuando rinden
mejor en las tareas médicas estándar~\cite{gu2024medvh}. La alucinación es especialmente costosa en medicina,
porque el modelo describe con la misma fluidez los hallazgos que la imagen contiene y los que no, y un
hallazgo inventado puede cambiar una conducta clínica.

La falla contraria se llama aquí \emph{hedging}. El modelo redacta descripciones genéricas, compatibles con
casi cualquier lesión, y evita comprometer un diagnóstico, por ejemplo enumerando varias categorías posibles
sin inclinarse por ninguna. Ese texto no contiene errores visibles y tampoco informa nada sobre la lesión
concreta. Expresar incertidumbre no es en sí una falla: se convierte en \emph{hedging} cuando es
indiscriminada, es decir, cuando no depende de cuán incierto sea el caso. Se ha argumentado que los esquemas
de evaluación que premian acertar y no penalizan adivinar incentivan la alucinación en lugar de la expresión
de incertidumbre~\cite{kalai2025hallucinate}. Enumerar diferenciales cuando una señal de confianza lo
justifica (Sección~\ref{sec:marco_tres_salidas}) es lo opuesto al \emph{hedging}.

\begin{table}[H]
    \centering
    \caption{Fallas de un reporte generado y su consecuencia para el lector, de mayor a menor gravedad.}
    \label{tab:marco_fallas}
    \small
    \begin{tabularx}{\textwidth}{@{}>{\raggedright\arraybackslash}p{3.0cm}
        >{\raggedright\arraybackslash\hspace{0pt}}X >{\raggedright\arraybackslash\hspace{0pt}}X@{}}
        \toprule
        \textbf{Falla} & \textbf{Qué hace el texto} & \textbf{Consecuencia para el lector} \\
        \midrule
        Diagnóstico incorrecto & Afirma una categoría distinta de la real
            & Recibe un error que la fluidez del texto hace parecer confiable \\
        \filasep
        Alucinación & Describe hallazgos que la imagen no contiene
            & Puede cambiar una conducta y no se distingue por la forma del texto \\
        \filasep
        Contradicción & Nombra dos o más diagnósticos incompatibles como conclusión
            & No sabe cuál de ellos sostiene el reporte \\
        \filasep
        \emph{Hedging} & Redacta un texto genérico y evita comprometer un diagnóstico
            & No recibe información y no sabe si el modelo duda o no \\
        \filasep
        Repetición & Repite frases en bucle~\cite{holtzman2020curious}
            & Recibe un texto degenerado, sin información útil \\
        \filasep
        Texto insuficiente & Una o dos frases sin descripción de la lesión
            & No recibe la descripción que el reporte debe contener \\
        \bottomrule
    \end{tabularx}
\end{table}

Estas dos fallas no son las únicas. Al leer los reportes uno por uno, cada uno se puede asignar a su falla
principal (Tabla~\ref{tab:marco_fallas}). Afirmar un diagnóstico incorrecto es la más grave, porque el lector
recibe un error con la forma de un acierto. La contradicción, que nombra dos diagnósticos incompatibles en el
mismo texto, y el \emph{hedging} dejan la decisión en el lector, y la repetición y el texto insuficiente no
aportan información. La alucinación puede acompañar a cualquiera de ellas, incluso a un diagnóstico correcto.
La taxonomía organiza el análisis de los reportes generados del Capítulo~\ref{cp:res_gen}.

Las métricas de solapamiento con un texto de referencia (Sección~\ref{sec:marco_eval_texto}) miden un
parecido superficial y pueden no detectar si el reporte transmite los hallazgos clínicos
correctos~\cite{hartsock2024vision}, por lo que la fidelidad
diagnóstica y el contenido clínico se miden por separado. Una manera de reducir estas fallas es darle al
modelo, por la vía de la Ecuación~\ref{eq:marco_rep_condicionada}, información que no obtiene de la imagen
sola, y la recuperación es la forma más estudiada de obtenerla.

\subsection{La recuperación como fuente de evidencia}
\label{sec:marco_rep_recuperacion}

\phantomsection\label{sec:marco_rag}%
La \gls{rag}, de 2020, agrega al generador un paso de búsqueda en una base externa de documentos. Tiene tres
partes: una base de documentos, un recuperador que selecciona los más relevantes para cada entrada y un
generador que produce la salida condicionado en la entrada y en lo recuperado. En la formulación original, en su variante
RAG-Sequence, un codificador proyecta cada documento $z$ de la base a un embedding $\mathbf{d}(z)$ y otro proyecta la consulta
$x$ a $\mathbf{q}(x)$. El recuperador asigna a cada documento una probabilidad
$p_\eta(z \mid x) \propto \exp\!\left(\mathbf{d}(z)^{\top}\mathbf{q}(x)\right)$, y el generador produce la
salida $y$ con cada uno de los $K$ documentos más probables y combina los resultados pesándolos por esa
probabilidad,
\begin{equation}
    p(y \mid x) \approx \sum_{z \in \text{top-}K} p_\eta(z \mid x)\, p_\theta(y \mid x, z).
    \label{eq:marco_rag}
\end{equation}
El codificador de la consulta se ajusta junto con el generador, mientras que el de documentos y el índice
quedan fijos, y la base se puede cambiar o ampliar sin reentrenar~\cite{lewis2020retrievalaugmented}. En
radiología, la recuperación multimodal parte de la imagen y recupera reportes de casos parecidos, y entrenar el
recuperador para que prefiera reportes con los mismos hallazgos clínicos mejora la corrección de lo
generado~\cite{sun2025factaware}.

\phantomsection\label{sec:marco_recuperacion_imagenes}%
Con una base de casos etiquetados, lo recuperable son imágenes. Los $k$ vecinos de la
Ecuación~\ref{eq:marco_vecinos} son los casos más parecidos a la consulta en el espacio de un codificador
visual, y su información, como su diagnóstico o su localización, es \textbf{evidencia visual}: dice qué se vio
que se parezca a la consulta. Aplicada a un modelo de lenguaje, la búsqueda de vecinos mejora la predicción sin
entrenamiento adicional, sobre todo en los patrones poco frecuentes, porque el modelo accede a casos concretos
que no tiene que memorizar~\cite{khandelwal2020generalization}. Un antecedente en descripción de imágenes
recupera las descripciones de imágenes parecidas y las entrega como \emph{prompt} a un modelo de lenguaje, sin
entrenar el recuperador~\cite{ramos2023smallcap}.

Esta segunda forma es un caso de la Ecuación~\ref{eq:marco_rep_condicionada} con
$c(x) = \psi\bigl(\mathcal{N}_k(x)\bigr)$, donde $\psi$ serializa como texto la información de los vecinos, y
difiere de la formulación de \textcite{lewis2020retrievalaugmented} en tres aspectos.
\begin{enumerate}
    \item El recuperador puede quedar congelado, sin entrenarse junto con el generador.
    \item No se marginaliza: el generador recibe todo el contexto a la vez y produce una única salida.
    \item Lo recuperado son etiquetas y atributos de imágenes, de la misma modalidad que la consulta, y no
    documentos de texto.
\end{enumerate}

\phantomsection\label{sec:marco_integracion}%
Recuperar documentos, recuperar imágenes y medir la concordancia entre predicciones de la
Sección~\ref{sec:marco_senales_confianza} no son equivalentes (Tabla~\ref{tab:marco_mecanismos}). Los dos
primeros aportan evidencia al texto y atienden el segundo eje de la Sección~\ref{sec:marco_rep_ejes}; la
concordancia no aporta evidencia, pero decide cuánta afirmación permite la que hay y atiende el tercero.
Integrados, la evidencia recuperada condiciona el texto y la señal de confianza elige si el texto afirma,
enumera diferenciales o no se redacta.

\begin{table}[H]
    \centering
    \caption{Mecanismos que agregan información al generador: qué aportan, su modalidad, qué ayudan a resolver
    y qué dejan abierto.}
    \label{tab:marco_mecanismos}
    \small
    \begin{tabularx}{\textwidth}{@{}>{\raggedright\arraybackslash}p{2.7cm}
        >{\raggedright\arraybackslash\hspace{0pt}}X >{\raggedright\arraybackslash\hspace{0pt}}X
        >{\raggedright\arraybackslash\hspace{0pt}}X >{\raggedright\arraybackslash\hspace{0pt}}X@{}}
        \toprule
        \textbf{Mecanismo} & \textbf{Qué recupera o calcula} & \textbf{Modalidad} & \textbf{Ayuda a resolver} & \textbf{No resuelve} \\
        \midrule
        Recuperar documentos (\gls{rag}) & Textos de una base, seleccionados por un recuperador
            & Consulta y documentos; los documentos son texto
            & Descripción y diagnóstico con conocimiento textual
            & Requiere una base de documentos pertinentes y no dice cuánto confiar en ellos \\
        \filasep
        Recuperar imágenes similares & Etiquetas y atributos de los $k$ casos más cercanos
            & Imagen contra imagen
            & Diagnóstico, como evidencia visual
            & No aporta redacción, y los vecinos existen aunque no se parezcan a la consulta \\
        \filasep
        Concordancia entre predicciones & Si dos predictores coinciden sobre la misma imagen
            & Predicciones
            & Incertidumbre, es decir, cuánta afirmación se permite
            & No agrega información al texto ni corrige un error que los dos predictores comparten \\
        \bottomrule
    \end{tabularx}
\end{table}

Clasificadores, generadores y señales de confianza producen así salidas de tres tipos: categorías, textos y
decisiones de abstención. Cada una exige sus propias métricas.

\section{Evaluación de modelos}
\label{sec:marco_evaluacion}

Las métricas se agrupan según lo que evalúan: la categoría que asigna un clasificador, el texto que genera un
modelo y las propiedades que no dependen del tipo de salida.

\subsection{Evaluación de la clasificación}
\label{sec:marco_eval_clasif}

El resumen más completo de un clasificador es la \textbf{matriz de confusión} $M \in \mathbb{N}^{C \times C}$,
en la que $M_{ij}$ cuenta las imágenes de la clase real $i$ que el modelo predice como $j$, y de la que salen
todas las demás métricas. En la tarea binaria, con una clase de interés (la positiva) frente al resto, tiene
cuatro celdas (Figura~\ref{fig:marco_matriz_roc}a): los verdaderos positivos ($\mathrm{TP}$), los falsos
negativos ($\mathrm{FN}$), los falsos positivos ($\mathrm{FP}$) y los verdaderos negativos ($\mathrm{TN}$). Cada
fila se lee contra la clase real y cada columna contra la predicha. Las métricas que este trabajo usa son las
siguientes.

% Figura: matriz de confusión binaria y curva ROC ilustrativas, con datos sintéticos (no son resultados del
% trabajo). Se usa en sec:marco_eval_clasif, al presentar la matriz de confusión y las métricas que salen de ella.

\begin{figure}[H]
    \centering
    \begin{tikzpicture}[x=1mm, y=1mm]

        % (a) Matriz de confusión binaria.
        \begin{scope}
            \foreach \x/\y/\txt/\col in {0/20/{TP}/figverde, 20/20/{FN}/figbermellon,
                                         0/0/{FP}/figbermellon, 20/0/{TN}/figverde} {
                \fill[\col!15] (\x, \y) rectangle ++(20, 20);
                \draw[figtrazo] (\x, \y) rectangle ++(20, 20);
                \node[font=\small\bfseries] at (\x + 10, \y + 10) {\txt};
            }
            \node[font=\footnotesize] at (10, 43) {positiva};
            \node[font=\footnotesize] at (30, 43) {negativa};
            \node[font=\small] at (20, 49) {clase predicha};
            \node[font=\footnotesize, anchor=east] at (-1.5, 30) {positiva};
            \node[font=\footnotesize, anchor=east] at (-1.5, 10) {negativa};
            \node[font=\small, rotate=90] at (-19, 20) {clase real};

            % Lectura por fila (sensibilidad, especificidad) y por columna (precisión).
            \draw[figflecha] (41, 30) -- (47, 30)
                node[right, font=\footnotesize, align=left] {sensibilidad\\$\mathrm{TP}/(\mathrm{TP}+\mathrm{FN})$};
            \draw[figflecha] (41, 10) -- (47, 10)
                node[right, font=\footnotesize, align=left] {especificidad\\$\mathrm{TN}/(\mathrm{TN}+\mathrm{FP})$};
            \draw[figflecha] (10, -1) -- (10, -7)
                node[below, font=\footnotesize, align=center] {precisión\\$\mathrm{TP}/(\mathrm{TP}+\mathrm{FP})$};
            \node[font=\small] at (20, -24) {(a) Matriz de confusión};
        \end{scope}

        % (b) Curva ROC.
        \begin{scope}[xshift=95mm]
            \fill[figazul!15] (0, 0) -- (0, 22) .. controls (4, 35) and (12, 39) .. (40, 40) -- (40, 0) -- cycle;
            \draw[figtrazo] (0, 0) rectangle (40, 40);
            \draw[figtrazo, dashed] (0, 0) -- (40, 40);
            \draw[figazul!80!black, thick] (0, 0) -- (0, 22) .. controls (4, 35) and (12, 39) .. (40, 40);
            \fill[fignaranja!80!black] (6.2, 32.4) circle (1.1);
            \node[figetiqueta, anchor=west] at (8, 28) {umbral $\tau$};
            \node[font=\footnotesize] at (27, 14) {AUC};
            \node[font=\footnotesize, rotate=38] at (22, 18.5) {azar};
            \foreach \t in {0, 0.5, 1} {
                \node[below, font=\footnotesize] at (\t*40, 0) {\t};
                \node[left, font=\footnotesize] at (0, \t*40) {\t};
            }
            \node[font=\small] at (20, -7) {FPR $= 1 -{}$especificidad};
            \node[font=\small, rotate=90] at (-8, 20) {TPR $={}$sensibilidad};
            \node[font=\small] at (20, -24) {(b) Curva ROC};
        \end{scope}

    \end{tikzpicture}
    \caption{Matriz de confusión binaria con las métricas que se leen en ella, y curva ROC con su área.}
    \label{fig:marco_matriz_roc}
\end{figure}


\begin{itemize}
    \item \textbf{Exactitud.} Es la fracción de imágenes bien clasificadas, $\sum_c M_{cc} / n$. Es engañosa
    cuando las clases están desbalanceadas (Sección~\ref{sec:marco_desbalance}): si una clase reúne dos tercios
    de las imágenes, un modelo que responda siempre esa clase alcanza una exactitud cercana a 0.67.

    \item \textbf{Sensibilidad macro.} La sensibilidad de la clase $c$, $R_c = M_{cc} / \sum_{j} M_{cj}$, es la
    fracción de sus imágenes que el modelo acierta, y se lee por fila. La sensibilidad macro promedia esas
    fracciones para que cada clase pese lo mismo, y corrige así el defecto de la exactitud,
    \begin{equation}
        \text{sensibilidad macro} = \frac{1}{C}\sum_{c=1}^{C} R_c.
        \label{eq:marco_acc_bal}
    \end{equation}
    En el caso binario se la llama exactitud balanceada~\cite{brodersen2010balanced}, y para varias clases es el
    promedio de las sensibilidades por clase~\cite{sokolova2009systematic}.

    \item \textbf{Precisión y F1.} La sensibilidad no dice cuántas de las predicciones de una clase son
    correctas. Esa es la precisión (\emph{precision}) de la clase $c$, $P_c$, que se lee por columna. El F1 las
    combina mediante su media armónica, y su promedio sobre las clases se puede hacer de dos maneras,
    \begin{equation}
        F1_c = \frac{2\,P_c R_c}{P_c + R_c}, \qquad
        F1_{\text{macro}} = \frac{1}{C}\sum_{c=1}^{C} F1_c, \qquad
        F1_{\text{ponderado}} = \sum_{c=1}^{C} \frac{n_c}{n}\, F1_c,
        \label{eq:marco_f1}
    \end{equation}
    donde $n_c$ es el número de imágenes de la clase $c$ y $n$ el total. El F1 ponderado queda dominado por la
    clase más frecuente, y el F1 macro da a las $C$ clases el mismo peso.

    \item \textbf{Sensibilidad y especificidad de una clase.} Cuando interesa una sola condición, la salida
    multiclase se colapsa en una binaria: la clase de interés, o un grupo de clases como las malignas, frente a
    las demás, con la probabilidad de la clase o la suma de las del grupo como puntaje. A esa tarea derivada se
    la llama detección, y se evalúa con la sensibilidad $R_c$ y con la especificidad, la fracción de las imágenes
    de las otras clases que el modelo no predice como $c$,
    \begin{equation}
        \mathrm{Esp}_c = \frac{\sum_{i \neq c}\sum_{j \neq c} M_{ij}}{\sum_{i \neq c}\sum_{j} M_{ij}}.
        \label{eq:marco_sens_esp}
    \end{equation}
    Las dos se mueven en sentidos opuestos cuando cambia el umbral de decisión~\cite{fawcett2006introduction}.

    \item \textbf{Curva ROC y \gls{auc}.} El modelo asigna a cada imagen un puntaje $s$ y la declara positiva si
    $s$ supera un umbral $\tau$. Cada umbral produce una tasa de verdaderos positivos, que es la sensibilidad, y
    una de falsos positivos, que es uno menos la especificidad,
    \begin{equation}
        \mathrm{TPR}(\tau) = \frac{\mathrm{TP}}{\mathrm{TP} + \mathrm{FN}}, \qquad
        \mathrm{FPR}(\tau) = \frac{\mathrm{FP}}{\mathrm{FP} + \mathrm{TN}}.
        \label{eq:marco_tpr_fpr}
    \end{equation}
    La curva ROC recorre todos los umbrales y dibuja la TPR contra la FPR (Figura~\ref{fig:marco_matriz_roc}b),
    y el \gls{auc} es el área bajo esa curva. No depende de un umbral, vale 0.5 para un clasificador al azar y 1
    para uno que separa perfectamente las dos clases, y es la probabilidad de que el modelo asigne más puntaje a
    un positivo elegido al azar que a un negativo elegido al azar,
    \begin{equation}
        \mathrm{AUC} = P\!\left(s(x^{+}) > s(x^{-})\right).
        \label{eq:marco_auc}
    \end{equation}
    Se calcula con la regla del trapecio. En el caso multiclase se construye una curva por clase, que trata esa
    clase como positiva y al resto como negativa~\cite{fawcett2006introduction}. Fawcett pondera las $C$ áreas por
    la prevalencia de cada clase; aquí se usa el promedio simple, que da a todas el mismo peso,
    $\mathrm{AUC}_{\text{macro}} = \frac{1}{C}\sum_{c} \mathrm{AUC}_c$.
\end{itemize}

Un reporte, en cambio, es un texto, y compararlo con otro exige métricas distintas.

\subsection{Evaluación de la generación}
\label{sec:marco_eval_texto}

Una misma imagen admite muchas descripciones válidas, y un texto puede parecerse a la referencia sin ser
correcto. Las métricas se agrupan en tres familias (Tabla~\ref{tab:marco_familias_texto}): las textuales
comparan el texto generado con uno de referencia, las diagnósticas comprueban si nombra el diagnóstico
correcto y las descriptivas cuentan cuánto contenido clínico menciona. Estas dos últimas no necesitan texto
de referencia.

\begin{table}[H]
    \centering
    \caption{Familias de métricas de generación de reportes: qué comparan, qué miden y qué no ven.}
    \label{tab:marco_familias_texto}
    \small
    \begin{tabularx}{\textwidth}{@{}>{\raggedright\arraybackslash}p{2.3cm}
        >{\raggedright\arraybackslash\hspace{0pt}}X >{\raggedright\arraybackslash\hspace{0pt}}X
        >{\raggedright\arraybackslash\hspace{0pt}}X >{\raggedright\arraybackslash\hspace{0pt}}X@{}}
        \toprule
        \textbf{Familia} & \textbf{Métricas} & \textbf{Compara con} & \textbf{Mide} & \textbf{No ve} \\
        \midrule
        Textuales & BLEU, ROUGE-L, METEOR, CIDEr, BERTScore & Un texto de referencia
            & Parecido léxico o de significado con esa referencia
            & Si lo descrito está en la imagen \\
        \filasep
        Diagnósticas & Consistencia diagnóstica & La etiqueta de la imagen
            & Si el texto nombra la categoría correcta
            & Si la descripción es fiel a la imagen \\
        \filasep
        Descriptivas & Cobertura de entidades clínicas & Un vocabulario cerrado
            & Cuánto contenido clínico menciona el texto
            & Si lo mencionado es cierto \\
        \bottomrule
    \end{tabularx}
\end{table}

Las métricas textuales comparan el texto generado $g$ con un texto de referencia $r$ escrito para la misma
imagen. Las cuatro primeras cuentan coincidencias de $n$-gramas, secuencias de $n$ palabras consecutivas, y la
quinta compara significados.

\begin{itemize}
    \item \textbf{BLEU.} Calcula, para cada $n$ entre 1 y 4, la precisión $p_n$ de los $n$-gramas de $g$, la
    fracción que también aparece en $r$, donde cada $n$-grama cuenta a lo sumo tantas veces como aparece en $r$
    (precisión recortada). Combina las cuatro con una media geométrica y multiplica el resultado
    por una penalización de brevedad, porque un texto muy corto puede tener precisión alta sin cubrir la
    referencia,
    \begin{equation}
        \mathrm{BLEU} = \mathrm{BP}\cdot\exp\!\left(\sum_{n=1}^{4} \tfrac{1}{4}\log p_n\right), \qquad
        \mathrm{BP} =
        \begin{cases}
            1 & \text{si } \ell_g > \ell_r,\\
            e^{\,1 - \ell_r/\ell_g} & \text{si } \ell_g \le \ell_r,
        \end{cases}
        \label{eq:marco_bleu}
    \end{equation}
    donde $\ell_g$ y $\ell_r$ son las longitudes en palabras de los dos textos~\cite{papineni2002bleu}.

    \item \textbf{ROUGE-L.} Se basa en la subsecuencia común más larga, $\mathrm{LCS}(g, r)$, la secuencia más
    larga de palabras que aparece en ambos textos en el mismo orden. Con ella define sensibilidad y precisión y
    las combina en una medida F,
    \begin{equation}
        R_{\mathrm{lcs}} = \frac{\mathrm{LCS}(g, r)}{\ell_r}, \qquad
        P_{\mathrm{lcs}} = \frac{\mathrm{LCS}(g, r)}{\ell_g}, \qquad
        \text{ROUGE-L} = \frac{(1+\beta^2)\,R_{\mathrm{lcs}}\,P_{\mathrm{lcs}}}{R_{\mathrm{lcs}} + \beta^2 P_{\mathrm{lcs}}},
        \label{eq:marco_rouge}
    \end{equation}
    donde $\beta$ fija cuánto pesa la sensibilidad frente a la precisión~\cite{lin2004rouge}. Con $\beta = 1$,
    la versión que usan las bibliotecas habituales, queda la media armónica de las dos.

    \item \textbf{METEOR.} Alinea las palabras de los dos textos, primero las idénticas, después las que
    comparten la raíz y por último los sinónimos. Sobre las $m$ palabras alineadas calcula la precisión $P$ y la
    sensibilidad $R$ y las combina dando más peso a la segunda. Luego agrupa las palabras alineadas en la menor
    cantidad posible de tramos contiguos, $ch$, y penaliza que las coincidencias estén fragmentadas,
    \begin{equation}
        F_{\text{mean}} = \frac{10\,P R}{R + 9P}, \qquad
        \mathrm{Pen} = 0.5\left(\frac{ch}{m}\right)^{3}, \qquad
        \mathrm{METEOR} = F_{\text{mean}}\,(1 - \mathrm{Pen}).
        \label{eq:marco_meteor}
    \end{equation}
    La penalización es casi nula si las coincidencias forman un único tramo y reduce el puntaje hasta la mitad
    si ninguna palabra alineada es contigua a otra~\cite{banerjee2005meteor}.

    \item \textbf{CIDEr.} Se diseñó para describir imágenes con varias referencias por imagen. Representa cada
    texto con un vector $\mathbf{t}^n(\cdot)$ con un peso TF-IDF por $n$-grama, que crece con su frecuencia en
    ese texto y decrece con la cantidad de imágenes en cuyas descripciones aparece, de modo que las expresiones
    comunes pesan poco. Para cada $n$ promedia la similitud coseno con las $J$ referencias $r_1, \dots, r_J$, y
    después promedia sobre $n$~\cite{vedantam2015cider},
    \begin{equation}
        \mathrm{CIDEr}_n = \frac{1}{J}\sum_{j=1}^{J}
        \frac{\mathbf{t}^n(g)^{\top}\mathbf{t}^n(r_j)}{\lVert\mathbf{t}^n(g)\rVert\,\lVert\mathbf{t}^n(r_j)\rVert},
        \qquad
        \mathrm{CIDEr} = \sum_{n=1}^{4} \tfrac{1}{4}\,\mathrm{CIDEr}_n.
        \label{eq:marco_cider}
    \end{equation}

    \item \textbf{BERTScore.} Las cuatro anteriores se apoyan en la coincidencia de palabras, con raíces y
    sinónimos en METEOR, y una paráfrasis correcta puede puntuar bajo. BERTScore compara en cambio los \emph{embeddings} contextuales de cada token, obtenidos con un
    modelo de solo codificador como BERT (Sección~\ref{sec:marco_llm_grandes}). Cada token de un texto se
    empareja con el más parecido del otro por similitud coseno, y las similitudes se promedian,
    \begin{equation}
        R_{\mathrm{BERT}} = \frac{1}{|r|}\sum_{r_i \in r}\,\max_{g_j \in g}\, \mathbf{r}_i^{\top}\mathbf{g}_j,
        \qquad
        P_{\mathrm{BERT}} = \frac{1}{|g|}\sum_{g_j \in g}\,\max_{r_i \in r}\, \mathbf{r}_i^{\top}\mathbf{g}_j,
        \label{eq:marco_bertscore}
    \end{equation}
    donde los \emph{embeddings} $\mathbf{r}_i$ y $\mathbf{g}_j$ están normalizados, y el puntaje final
    $F_{\mathrm{BERT}}$ es la media armónica de los dos, como el F1 de la Ecuación~\ref{eq:marco_f1}. Como dos
    sinónimos tienen \emph{embeddings} cercanos, la métrica tolera reformulaciones~\cite{zhang2020bertscore}.
\end{itemize}

Las métricas con referencia suponen que la referencia es la descripción correcta de la imagen, y en
dermatología ese supuesto no está garantizado: cuando varios dermatólogos marcan por separado los mismos
criterios dermatoscópicos, la concordancia es en su mayoría de pobre a regular, y moderada solo en unos pocos
criterios~\cite{carrera2016validity}. Una descripción
generada puede diferir de la referencia sin estar equivocada. Las métricas diagnósticas y descriptivas evitan
ese supuesto porque comprueban lo que el texto afirma.

\begin{itemize}
    \item \textbf{Consistencia diagnóstica.} Mide la fidelidad diagnóstica de la
    Sección~\ref{sec:marco_rep_ejes}: la fracción de los textos que nombran la categoría real de la imagen o un
    sinónimo reconocido,
    \begin{equation}
        \mathrm{FD} = \frac{1}{n}\sum_{i=1}^{n} \mathbf{1}\bigl[\mathrm{dx}(g_i) = y_i\bigr],
        \label{eq:marco_fd}
    \end{equation}
    donde $g_i$ es el texto generado para la imagen $i$, $y_i$ su categoría real y $\mathrm{dx}(\cdot)$ extrae la
    categoría que el texto nombra.

    \item \textbf{Cobertura de entidades clínicas.} Se apoya en el vocabulario cerrado de la
    Sección~\ref{sec:marco_rep_tarea}. Es la proporción de las categorías del vocabulario, como color, patrón o
    estructura, que el texto menciona,
    \begin{equation}
        \mathrm{cob}(g) = \frac{\bigl|\{k : \exists\, e \in V_k,\ e \text{ aparece en } g\}\bigr|}{K},
        \label{eq:marco_cobertura_ent}
    \end{equation}
    donde $V_k$ es el conjunto de términos de la categoría $k$ y $K$ el número de categorías. Mide cuánto
    contenido clínico describe el texto y no si lo que describe es cierto.
\end{itemize}

Ninguna de estas familias comprueba que lo descrito esté en la imagen, que es lo que define a una
alucinación~\cite{hartsock2024vision}. Detectarla de la Sección~\ref{sec:marco_rep_fallas} requiere que un especialista lea el reporte contra la
imagen y lo puntúe según una rúbrica, es decir, según criterios fijados de antemano. Las propiedades que
siguen se aplican a ambos tipos de salida.

\subsection{Evaluación transversal}
\label{sec:marco_eval_transversal}

\setcounter{secnumdepth}{3}
\subsubsection{Calibración}
\label{sec:marco_eval_calib}

La calibración se definió en la Sección~\ref{sec:marco_senales_confianza} como la coincidencia entre las
probabilidades de un modelo y sus frecuencias de acierto. El diagrama de confiabilidad la hace visible
(Figura~\ref{fig:marco_confiabilidad}): las predicciones se agrupan en $M$ intervalos según su confianza, y
para cada intervalo $B_m$ se compara la exactitud $\mathrm{acc}(B_m)$ con la confianza media
$\mathrm{conf}(B_m)$. En un modelo calibrado las barras tocan la diagonal, y la brecha entre ambas es la que
promedia el ECE. El \textbf{\gls{ece}} la pesa por la cantidad de predicciones de cada intervalo, y el
\textbf{error de calibración máximo} (MCE) toma la peor,
\begin{equation}
    \mathrm{ECE} = \sum_{m=1}^{M} \frac{|B_m|}{n}\,\bigl|\mathrm{acc}(B_m) - \mathrm{conf}(B_m)\bigr|,
    \qquad
    \mathrm{MCE} = \max_{m}\,\bigl|\mathrm{acc}(B_m) - \mathrm{conf}(B_m)\bigr|.
    \label{eq:marco_ece}
\end{equation}
El ECE resume la calibración típica del modelo, y el MCE importa cuando un error de confianza en un intervalo
concreto tiene consecuencias graves~\cite{guo2017calibration}.

El \textbf{puntaje de Brier} mide la calibración sin agrupar en intervalos. Es el error cuadrático medio entre
el vector de probabilidades y la codificación de la clase real,
\begin{equation}
    \mathrm{BS} = \frac{1}{n}\sum_{i=1}^{n}\sum_{c=1}^{C} \left(p_{ic} - y_{ic}\right)^2,
    \label{eq:marco_brier}
\end{equation}
donde $p_{ic}$ es la probabilidad que el modelo asigna a la clase $c$ en la imagen $i$, e $y_{ic}$ vale 1 si
$c$ es la clase real de esa imagen y 0 si no~\cite{brier1950verification}. Baja tanto cuando las
probabilidades se acercan a las frecuencias reales como cuando el modelo concentra más probabilidad en la
clase correcta, y no mide solo la calibración.

% Figura: diagrama de confiabilidad ilustrativo, con datos sintéticos (no son resultados del trabajo).
% Se usa en sec:marco_eval_calib, después de la definición de ECE y MCE.

\begin{figure}[!tbp]
    \centering
    \begin{tikzpicture}[x=6.5mm, y=55mm]

        % Ejes.
        \draw[figtrazo] (0, 0) rectangle (10, 1);
        \foreach \t in {0, 0.2, 0.4, 0.6, 0.8, 1.0} {
            \draw[figtrazo] (\t*10, 0) -- ++(0, -0.02) node[below, font=\footnotesize] {\t};
            \draw[figtrazo] (0, \t) -- ++(-0.12, 0) node[left, font=\footnotesize] {\t};
        }
        \node[font=\small] at (5, -0.14) {confianza de la predicción};
        \node[font=\small, rotate=90] at (-1.35, 0.5) {exactitud};

        % Barras: exactitud por intervalo (sintética, modelo sobreconfiado).
        \foreach \b/\acc in {0/0.04, 1/0.12, 2/0.21, 3/0.28, 4/0.36, 5/0.44, 6/0.53, 7/0.61, 8/0.70, 9/0.82} {
            \pgfmathsetmacro{\medio}{(\b + 0.5)/10}
            \fill[figazul!35] (\b + 0.05, 0) rectangle (\b + 0.95, \acc);
            \draw[figazul!80!black] (\b + 0.05, 0) rectangle (\b + 0.95, \acc);
            % Brecha hasta la diagonal.
            \fill[figbermellon!30, opacity=0.8] (\b + 0.05, \acc) rectangle (\b + 0.95, \medio);
            \draw[figbermellon!80!black, densely dashed] (\b + 0.05, \acc) rectangle (\b + 0.95, \medio);
        }

        % Diagonal de calibración perfecta.
        \draw[figtrazo, thick, dashed] (0, 0) -- (10, 1);

        % Leyenda.
        \fill[figazul!35, draw=figazul!80!black] (0.5, 0.90) rectangle (1.1, 0.95);
        \node[anchor=west, font=\footnotesize] at (1.2, 0.925) {exactitud del intervalo};
        \fill[figbermellon!30, draw=figbermellon!80!black, densely dashed] (0.5, 0.81) rectangle (1.1, 0.86);
        \node[anchor=west, font=\footnotesize] at (1.2, 0.835) {brecha con la confianza};
        \draw[figtrazo, thick, dashed] (0.5, 0.745) -- (1.1, 0.745);
        \node[anchor=west, font=\footnotesize] at (1.2, 0.745) {calibración perfecta};

    \end{tikzpicture}
    \caption{Diagrama de confiabilidad de un modelo sobreconfiado, con datos sintéticos ilustrativos.}
    \label{fig:marco_confiabilidad}
\end{figure}


La calibración se puede corregir después del entrenamiento con un conjunto de validación. El \emph{temperature
scaling} divide todos los logits por un único escalar $T>0$ antes de la softmax, como en la
Ecuación~\ref{eq:marco_temperatura}, y ajusta $T$ para minimizar la entropía cruzada en validación. Como no
cambia el orden de los logits, la clase predicha y la exactitud quedan iguales. El \emph{vector scaling}
aprende en cambio un factor y un sesgo para cada clase; corrige más, pero puede cambiar la clase predicha y
con ello la exactitud~\cite{guo2017calibration}. Ni el ECE ni el puntaje de Brier dicen si la confianza sirve
para abstenerse, propiedad que mide la curva riesgo--cobertura de la Sección~\ref{sec:marco_pred_selectiva}.

\subsubsection{Inferencia estadística}
\label{sec:marco_eval_estadistica}

El conjunto de test es una muestra de las imágenes posibles, y cuando dos modelos difieren en una o dos
centésimas de F1 macro, esa diferencia puede deberse a qué imágenes cayeron en el test y no a los modelos.
Para afirmar que un modelo es mejor que otro hace falta estimar cuánto variaría la diferencia de una muestra
a otra.

El \textbf{\emph{bootstrap}} estima esa variación con los datos disponibles. Se construyen muchas muestras del
mismo tamaño que el test tomando imágenes al azar con reposición, se calcula la métrica en cada una, y la
dispersión de los valores obtenidos aproxima la que tendría la métrica entre muestras
independientes~\cite{efron1979bootstrap}. En la versión pareada, cada remuestreo toma las mismas imágenes para
los dos modelos y registra la diferencia entre ellos. Si el intervalo de percentiles de esas diferencias no
contiene el cero, la diferencia no se explica por la composición del test. Sirve para cualquier métrica
agregada, como el F1 macro o el AUC.

\begingroup La \textbf{prueba de McNemar} compara los aciertos de dos clasificadores sobre las mismas imágenes.
Solo usa los casos discordantes: $b$, los que el primero acierta y el segundo falla, y $c$, los del caso
inverso. Si los dos modelos fueran equivalentes, $b$ y $c$ deberían ser parecidos, y el
estadístico\displaywidowpenalty=10000
\begin{equation}
    \chi^2 = \frac{(b-c)^2}{b+c}
    \label{eq:marco_mcnemar}
\end{equation}
mide cuánto se separan; bajo la hipótesis de equivalencia sigue aproximadamente una distribución $\chi^2$ con
un grado de libertad~\cite{mcnemar1947note}. Cuando los casos discordantes son pocos se usa la versión
exacta: $b$ sigue una distribución binomial con $b+c$ ensayos y probabilidad $1/2$, y el valor $p$ se calcula
directamente de ella.\par\endgroup

\subsubsection{Robustez}
\label{sec:marco_eval_robustez}

\begingroup La robustez es cuánto cambia una métrica cuando la entrada se perturba de una manera que no altera su
contenido clínico, como el desenfoque, los cambios de brillo o de contraste, la compresión o el ruido. Sea $m$
una métrica calculada sobre el conjunto de test original y $m_s$ la misma métrica sobre el conjunto perturbado
con una perturbación de severidad $s$. La caída de desempeño, en valor absoluto o relativa a $m$,\displaywidowpenalty=10000
\begin{equation}
    \Delta m_s = m - m_s, \qquad
    \Delta m_s^{\text{rel}} = \frac{m - m_s}{m},
    \label{eq:marco_robustez}
\end{equation}
mide la sensibilidad a esa perturbación, y un modelo es más robusto cuanto menor es.\par\endgroup

Permite anticipar el comportamiento ante las variaciones de adquisición de la
Sección~\ref{sec:marco_generalizacion} sin otro conjunto de datos, degradando las imágenes del test de
maneras controladas~\cite{hendrycks2019benchmarking}. Las perturbaciones sintéticas reproducen solo una parte
de la variación real entre equipos y centros, de modo que una caída pequeña no garantiza robustez ante un
cambio de dominio.

\subsubsection{Interpretabilidad visual}
\label{sec:marco_eval_gradcam}
\setcounter{secnumdepth}{2}

Las métricas anteriores dicen cuánto acierta un clasificador pero no en qué se basa. \textbf{Grad-CAM}
(\emph{gradient-weighted class activation mapping}) produce un mapa de calor de las regiones de la imagen que
más influyeron en la predicción de una clase. Parte de los mapas de características $A^k$ de la última capa
convolucional, de $i \times j$ posiciones cada uno, y del puntaje $y^c$ de la clase $c$ antes de la softmax.
Mide cuánto importa cada mapa promediando el gradiente del puntaje respecto del mapa, suma los mapas pesados
por esa importancia y conserva solo los valores positivos,
\begin{equation}
    \alpha_k^c = \frac{1}{Z}\sum_{i}\sum_{j} \frac{\partial y^c}{\partial A^k_{ij}}, \qquad
    L^c_{\text{Grad-CAM}} = \mathrm{ReLU}\!\left(\sum_k \alpha_k^c A^k\right),
    \label{eq:marco_gradcam}
\end{equation}
donde $Z$ es la cantidad de posiciones del mapa. La ReLU descarta las regiones que bajan el puntaje de la
clase, que suelen corresponder a otras clases~\cite{selvaraju2017gradcam}. El mapa se reescala al tamaño de
la imagen y se superpone a ella. En el Swin, cuya última etapa también devuelve un mapa con disposición
espacial (Sección~\ref{sec:marco_swin}), el cálculo es igual.

El mapa tiene la resolución de la última capa, mucho menor que la de la imagen, y por eso marca regiones
gruesas~\cite{selvaraju2017gradcam}. Que se vea razonable tampoco basta: algunos métodos de saliencia
producen mapas casi iguales con los pesos del modelo aleatorizados, y la inspección visual por sí sola no
distingue un método fiel de uno que no lo es~\cite{adebayo2018sanity}. El mapa muestra dónde se apoyó el
modelo y no si esa región justifica el diagnóstico, pero sirve para detectar un modelo que se apoya en algo
ajeno a la lesión, como una regla, un borde oscuro o un marcador en la piel. Esa clase de atajos lleva a la
pregunta de hasta dónde valen las métricas fuera de la distribución del test.

\section{Sesgo, robustez y generalización}
\label{sec:marco_generalizacion}

Las métricas de la Sección~\ref{sec:marco_evaluacion} suelen calcularse sobre un conjunto de test de la misma
fuente que el de entrenamiento, pero en uso real las imágenes llegan de otros equipos, centros y pacientes.
Ese cambio de distribución (\emph{distribution shift}) reduce el desempeño de forma sustancial en problemas de
todo tipo, y la caída persiste aun con los métodos diseñados para mitigarla. Se distinguen el cambio de
dominio, cuando las imágenes provienen de una fuente que el entrenamiento no vio, y el cambio de
subpoblación, cuando la fuente es la misma pero los grupos aparecen en otras
proporciones~\cite{koh2021wilds}.

En dermatología el cambio de distribución más visible es el de modalidad. La fotografía clínica registra la
lesión tal como se ve a simple vista, y la dermatoscopía la examina con aumento e iluminación controladas, con
lo que muestra estructuras que la fotografía no resuelve, como la red pigmentaria y los patrones
vasculares~\cite{dinnes2018dermoscopy,kittler2016standardization}. Los rasgos visuales tienen distribuciones
distintas en uno y otro registro, así que un modelo entrenado con una modalidad y usado con la otra enfrenta
un cambio de dominio.

Aun dentro de una misma modalidad, las condiciones de adquisición varían entre equipos y centros: la
iluminación, la distancia a la lesión, la resolución del sensor y la compresión de la imagen. Un modelo puede
apoyarse en rasgos que esas variaciones alteran, como la textura fina o el color exacto, y perder desempeño
ante imágenes de contenido clínico idéntico. Esa fragilidad se estima perturbando las imágenes de test de
manera controlada (Sección~\ref{sec:marco_eval_robustez})~\cite{hendrycks2019benchmarking}.

Un conjunto puede contener además rasgos que acompañan a una clase sin formar parte de la lesión, como las
marcas de tinta con que se señala una lesión antes de biopsiarla, las etiquetas adhesivas y los bordes oscuros
del dermatoscopio. Si en el conjunto esos rasgos aparecen con más frecuencia en los melanomas, un modelo puede
aprender a predecir melanoma cuando los ve y alcanzar una exactitud alta sin mirar la lesión, un atajo que se
ha documentado en conjuntos públicos de dermatoscopía~\cite{yan2023towards}. Los propios modelos contrastivos
médicos permiten buscar esos artefactos en un conjunto antes de entrenar sobre él~\cite{kim2024transparent}, y
un mapa de Grad-CAM que marca la tinta y no la lesión (Sección~\ref{sec:marco_eval_gradcam}) es el síntoma
visible del atajo.

El tono de piel es un caso de cambio de subpoblación. Cuando las pieles claras están mucho más representadas
que las oscuras en el entrenamiento, el modelo clasifica mejor los fototipos parecidos a los que
vio~\cite{groh2021neural}. Evaluados sobre imágenes con fototipos diversos y confirmación histopatológica,
modelos dermatológicos del estado del arte rindieron bastante peor que en sus conjuntos de origen, sobre todo
en pieles oscuras y en enfermedades poco comunes, y el ajuste fino con imágenes de fototipos diversos cerró la
brecha~\cite{daneshjou2022disparities}.

\begingroup La disparidad no desaparece cuando el modelo asiste a un médico en lugar de decidir solo. En un
experimento con 389 dermatólogos y 459 médicos de atención primaria, el apoyo de un sistema de aprendizaje
profundo mejoró la exactitud de los dos grupos, pero en los médicos de atención primaria amplió la diferencia
entre pieles claras y oscuras~\cite{groh2024learningaided}.\clubpenalty=10000 \par\endgroup

La población de destino tampoco es la de los conjuntos de entrenamiento. Una revisión sistemática de 272
estudios sobre detección de cáncer de piel con aprendizaje automático encontró que muchos usaron conjuntos de
Europa y Estados Unidos con fototipos mayormente claros, y que solo dos usaron datos de poblaciones de baja
prevalencia, no derivadas a un especialista; por eso los autores no recomiendan todavía adoptar estos modelos en atención
primaria~\cite{jones2022intelligence}. Para la población argentina existe desde 2023 un conjunto público propio,
creado para evaluar herramientas de este tipo en ella~\cite{riccilara2023dataset}.

Con qué datos y con qué protocolo se evaluó un trabajo decide entonces si su resultado se sostiene fuera de su
conjunto de origen, y ese es el criterio con el que la sección siguiente revisa la literatura.

\section{Estado del arte}
\label{sec:estado_arte}

Los trabajos más cercanos a esta tesina se agrupan en tres ejes: la clasificación de lesiones de piel, los
\glspl{vlm} para dermatología y la generación de reportes dermatológicos. La revisión se concentra en trabajos
publicados desde 2020, y los anteriores aparecen solo como hitos, con su fecha. Varios de 2025 y 2026 están
publicados solo en arXiv o medRxiv, sin revisión de pares, y en cada caso se lo indica.

\subsection{Diagnóstico automático de lesiones de piel}
\label{sec:estado_arte_diagnostico}

La clasificación de lesiones asigna a una imagen una entre varias categorías diagnósticas, sobre imágenes
dermatoscópicas (HAM10000 y los desafíos ISIC, conjuntos de referencia muy utilizados) o fotografías clínicas.
Importa cómo trata cada trabajo el desbalance de clases, con qué métrica elige el modelo, si mantiene las imágenes
de una misma lesión en una partición, si valida fuera de la fuente y si compara contra clínicos. La
Tabla~\ref{tab:marco_sota_clasificacion} resume los trabajos.

\begin{table}[!t]
    \centering
    \caption{Trabajos de clasificación de lesiones de piel. PCP: médicos de atención primaria; NP: enfermeros
    especializados (\emph{nurse practitioners}).}
    \label{tab:marco_sota_clasificacion}
    \footnotesize
    \setlength{\tabcolsep}{3pt}
    \begin{tabularx}{\textwidth}{@{}>{\raggedright\arraybackslash}p{1.9cm}
        c
        >{\raggedright\arraybackslash\hspace{0pt}}X >{\raggedright\arraybackslash\hspace{0pt}}X
        >{\raggedright\arraybackslash\hspace{0pt}}p{1.8cm} >{\raggedright\arraybackslash\hspace{0pt}}X
        >{\raggedright\arraybackslash\hspace{0pt}}p{1.75cm}@{}}
        \toprule
        \textbf{Trabajo} & \textbf{Año} & \textbf{Arquitectura} & \textbf{Datos} & \textbf{Métricas} & \textbf{Resultado principal} & \textbf{Expertos} \\
        \midrule
        \textcite{liu2020deep} & 2020 & Ensamble de 6 Inception-v4 con metadatos
            & 17\,777 casos de 17 sitios; 26 condiciones
            & Exactitud, sensibilidad promedio
            & Top-1 0.67; dermatólogos 0.63, PCP 0.45, NP 0.41
            & 18 lectores; 3 dermatólogos de referencia \\
        \filasep
        \textcite{gessert2020skin} & 2020 & Ensamble de CNN con metadatos
            & ISIC 2019; 5 particiones por lesión
            & Sensibilidad macro, AUC
            & Sin metadatos: 72.5\,\% en validación cruzada, 63.6\,\% en prueba
            & No \\
        \filasep
        \textcite{yao2022single} & 2022 & RegNetY única con pérdida propia
            & ISIC 2017, 2018, 2019 y 7-point
            & Sensibilidad macro y de melanoma, AUC
            & ISIC 2018: 0.875 y 0.819; ISIC 2019: 0.645
            & No \\
        \filasep
        \textcite{marchetti2023prospective} & 2023 & Ensamble de 18 modelos
            & 603 lesiones biopsiadas, 95 melanomas
            & Sensibilidad, especificidad, AUC
            & Sensibilidad 96.8\,\%; especificidad 37.4\,\%; AUC 0.857
            & 11 dermatólogos; impacto teórico \\
        \filasep
        \textcite{kreouzi2024learning} & 2024 & Cuatro CNN preentrenadas, entre ellas DenseNet121
            & DermNet; melanoma contra nevo
            & Exactitud, AUC, tiempo de inferencia, tamaño
            & Exactitud 92.30\,\% (DenseNet121); AUC 0.957 (ResNet50V2)
            & No \\
        \filasep
        \textcite{aboulmira2024skinhealthmate} & 2024 & Ensamble de EfficientNetB1 y B5 en una aplicación
            & HAM10000; 7 clases
            & Exactitud, F1 macro, AUC
            & Exactitud 93\,\%; F1 macro 0.81; sensibilidad de melanoma 0.49
            & No \\
        \filasep
        \textcite{khan2025ensemble} & 2025 & Ensamble de tres CNN, más LLM
            & Nevo contra carcinoma basocelular; 80/10/10 al azar
            & Exactitud, F1, pérdida logarítmica
            & Exactitud 97.67\,\% contra 93.83\,\% (ResNet50)
            & No \\
        \filasep
        \textcite{kim2025skinsavvy2} & 2025 & Swin o EfficientNetV2, más GPT
            & HAM10000 ampliado con imágenes externas
            & Exactitud, precisión, sensibilidad, F1
            & Exactitud balanceada 0.780; F1 promedio 0.793 (Swin)
            & 2 dermatólogos entre 20 evaluadores \\
        \bottomrule
    \end{tabularx}
\end{table}

Dos hitos fijaron la tarea como un terreno en el que el aprendizaje profundo compite con el especialista. En 2017
una \gls{cnn} igualó a dermatólogos certificados en dos distinciones entre lesiones malignas y benignas, carcinoma
queratinocítico contra queratosis seborreica y melanoma contra nevo~\cite{esteva2017dermatologistlevel}. En 2019,
los clasificadores del desafío ISIC 2018, entrenados con HAM10000, acertaron en su prueba en promedio más que 511 participantes humanos, y los tres mejores superaron también a los expertos con
más de diez años de experiencia~\cite{tschandl2019comparison}.

\textcite{liu2020deep} extendieron la tarea de la lesión aislada a las 26 condiciones de piel más frecuentes en
atención primaria, a partir de fotografías clínicas y de la historia del paciente. Un ensamble de seis redes
Inception-v4 fusiona las imágenes de cada caso con 45 variables de metadatos y se entrena con etiquetas blandas
construidas con los votos de un panel de dermatólogos. Los datos son casos de teledermatología de 17 sitios de
Estados Unidos, con partición temporal y separada por paciente, y la referencia es el diagnóstico de tres dermatólogos, sin confirmación
patológica. Se mide la exactitud y la sensibilidad promedio entre condiciones. Frente a 18 clínicos, en un
subconjunto enriquecido con condiciones raras, alcanzó una exactitud top-1 de 0.67 contra 0.63 de los dermatólogos y 0.45
de los médicos de atención primaria; entre los tumores malignos, en cambio, su sensibilidad top-1 fue 0.57 contra
0.71 (0.88 contra 0.89 entre los tres primeros diagnósticos). Los
autores reconocen que la partición temporal no reemplaza a un conjunto externo.

\textcite{gessert2020skin} describen el ensamble de EfficientNet, SENet154 y ResNeXt que ganó las dos tareas del
desafío ISIC 2019. Ponderan la entropía cruzada por la frecuencia inversa de cada clase y agregan una novena clase,
con imágenes externas, para la clase desconocida de la prueba. Parten de 25\,331 imágenes, más datos externos, en cinco
particiones que mantienen juntas las de una misma lesión. La métrica es la sensibilidad macro, y la evidencia es el primer puesto, sin
comparación con otros métodos. En la tarea sin metadatos llegan a 72.5\,\% en validación cruzada, pero a 63.6\,\% en la
prueba oficial, una brecha que atribuyen en parte a la clase desconocida; con metadatos, 74.2\,\% y 63.4\,\%. No hubo expertos.

\textcite{yao2022single} proponen en cambio un modelo único, de bajo costo, para conjuntos pequeños y
desbalanceados. De 41 \glspl{cnn} eligen una RegNetY y le suman una pérdida con un coeficiente por clase, mayor para
el melanoma, y un aprendizaje acumulativo que pasa del entrenamiento convencional al rebalanceo. Usan ISIC 2017, 2018,
2019 y 7-point, con la sensibilidad macro como métrica principal y la del melanoma aparte, y se comparan con los
ganadores de los desafíos. En ISIC 2019 alcanzan una sensibilidad macro de 0.645 contra 0.636 del ganador original del
desafío, y en ISIC 2018 muestran que subir el coeficiente del melanoma eleva su sensibilidad de 0.772 a 0.871 a
costa de la sensibilidad macro. No
hubo expertos. Reconocen que los ensambles rinden en general mejor que un modelo único. El artículo no describe una
partición de validación: la selección se hace sobre la prueba de ISIC 2018.

\textcite{marchetti2023prospective} llevaron fuera del banco de pruebas el ganador del desafío SIIM-ISIC 2020,
un ensamble de 18 modelos EfficientNet y ResNet; cada modelo se entrenó en cinco
particiones, de modo que para cada imagen se ejecutan y promedian 90 redes. Lo validaron de forma
prospectiva, con biopsia como referencia, en un estudio observacional con 603 lesiones, 95 de ellas melanomas. Con
el umbral fijado de antemano, la sensibilidad fue 96.8\,\% y la especificidad 37.4\,\%, y el \gls{auc} de 0.857 fue
similar al de la prueba del desafío y superior al de los dermatólogos sin ayuda (0.780). Con acceso al puntaje, el
\gls{auc} de 11 dermatólogos subió a 0.816, aunque el impacto es teórico porque no se modificó la atención de ningún
paciente. Los autores reconocen que fue un solo centro, con una muestra pequeña y sin lesiones biopsiadas de pacientes de
fototipo V o VI.

Los trabajos recientes sobre HAM10000 y conjuntos afines se comparan casi siempre entre arquitecturas, sin expertos
ni un protocolo común. \textcite{kreouzi2024learning} enfrentan cuatro \glspl{cnn} (DenseNet121, ResNet50V2,
NASNetMobile y MobileNetV2), preentrenadas en ImageNet y con la misma cabeza de clasificación, en la distinción binaria
entre melanoma y nevo melanocítico sobre imágenes dermatoscópicas de DermNet. Además de la exactitud y el \gls{auc}
miden el tiempo de inferencia y el tamaño del modelo, y evalúan la significancia de las diferencias con la prueba de
McNemar. DenseNet121 logra la mayor exactitud (92.30\,\%) y ResNet50V2 el mayor \gls{auc} (0.957), pero MobileNetV2
llega a 92.19\,\% con el menor tamaño (9.89 MB) y la inferencia más rápida (23.46 ms), lo que la perfila para
dispositivos de pocos recursos. No hubo expertos ni una referencia externa. Los autores reconocen que la tarea es
binaria, que no cubre nevos atípicos ni otras lesiones pigmentadas, que las imágenes se reducen a 224$\times$224
píxeles y que el modelo no incorpora metadatos del paciente.

\textcite{aboulmira2024skinhealthmate} integran en la aplicación web y móvil SkinHealthMate un ensamble que promedia
las probabilidades de EfficientNetB1 y EfficientNetB5, con preprocesamiento (filtro de mediana) y aumento de datos,
para clasificar las siete categorías de HAM10000. La aplicación permite a los dermatólogos cargar la imagen, validar
el diagnóstico y dar seguimiento al paciente. Reportan una exactitud de 93\,\% y un F1 macro de 0.81, y atribuyen la
diferencia entre ambos a las clases con pocas muestras: en melanoma, la precisión es 0.63 y la sensibilidad 0.49. No
comparan el ensamble con otros métodos ni con clínicos. Reconocen que el modelo, entrenado con HAM10000, no garantiza
un diagnóstico correcto en pieles oscuras, que la calidad de la imagen afecta el resultado y que las condiciones
raras pueden clasificarse mal, por lo que lo presentan como complemento del criterio médico.

\textcite{khan2025ensemble}, publicado solo en arXiv, sin revisión de pares, combinan un ensamble heterogéneo de
EfficientNetB3, ResNet50 y DenseNet121 con un modelo de lenguaje que redacta un reporte a partir de la predicción. El
clasificador distingue nevo de carcinoma basocelular en imágenes dermatoscópicas, con partición estratificada al azar
80/10/10, y decide por mayoría de votos; los casos en que los modelos discrepan se marcan para revisión de un
especialista. Se compara contra cada red por separado con exactitud, F1 y pérdida logarítmica: el ensamble alcanza una
exactitud de 97.67\,\% contra 93.83\,\% de ResNet50, el mejor modelo individual. El texto lo genera LLaMA-3 70B con un
\emph{prompt} que recibe la enfermedad, la confianza y si hubo consenso, y se valora con una rúbrica propia sin
comparar con reportes clínicos. No participaron expertos. Los autores reconocen que el marco es binario, que la
diversidad del conjunto limita el desempeño en distintos tonos de piel y que falta una validación en entorno real.

\textcite{kim2025skinsavvy2} reemplazan la DenseNet de su sistema anterior por un transformer Swin o una EfficientNetV2
y amplían HAM10000 con imágenes de revistas médicas, el archivo ISIC y DermNet NZ para atenuar el desbalance, sin
aumento de datos para no alterar el color ni la textura de la lesión. Las exactitudes por clase del modelo entrenado con
la versión ampliada superan las del entrenado con HAM10000 en casi todas las clases, salvo en la mayoritaria. Con Swin
se obtienen una exactitud balanceada de 0.780, una precisión promedio de 0.812, una sensibilidad promedio de 0.779 y un F1
promedio de 0.793, contra 0.744, 0.727 y 0.734 del sistema anterior en precisión, sensibilidad y F1; la
comparación con cuatro trabajos previos, entre ellos el de Esteva et al., da una exactitud entre 0.705 y 0.776.
Sobre la salida del clasificador un modelo GPT genera recomendaciones de cuidado personalizadas, que valoraron 20
participantes, dos de ellos dermatólogos; esa evaluación mide las recomendaciones y no el diagnóstico. Los autores
reconocen una brecha de dominio entre el conjunto curado y las imágenes tomadas en entornos no controlados.

Entre 2020 y 2023, la línea pasó de la lesión aislada a las condiciones de atención
primaria~\cite{liu2020deep}, corrigió el desbalance dentro de los desafíos ISIC~\cite{gessert2020skin,yao2022single}
y probó un ensamble en un estudio prospectivo~\cite{marchetti2023prospective}. Esos trabajos muestran que un
promedio global oculta la debilidad en los tumores malignos, que el melanoma exige su propia sensibilidad y que el
\gls{auc} se sostiene fuera del desafío, aunque con una especificidad baja y una mejora modesta sobre los
dermatólogos. Los recientes repiten ensambles y comparaciones entre redes, y
reportan exactitudes altas que el propio F1 macro de \textcite{aboulmira2024skinhealthmate} relativiza. Dos de ellos
agregan un texto generado por un modelo de lenguaje sobre la predicción~\cite{khan2025ensemble,kim2025skinsavvy2}.

Quedan cuatro cuestiones abiertas. Casi todos los trabajos usan redes convolucionales, solo
\textcite{kim2025skinsavvy2} incluyen un transformer, y no hay un protocolo común entre trabajos. Solo
\textcite{gessert2020skin}, por lesión, y \textcite{liu2020deep}, por paciente, describen particiones que mantienen
juntos los casos relacionados; \textcite{yao2022single} seleccionan sobre la prueba y \textcite{khan2025ensemble}
parten al azar. La validación fuera de la fuente es temporal o de un solo centro sin lesiones de fototipos V y
VI, los demás evalúan en particiones o pruebas de la misma fuente, y ninguno se desarrolla ni se valida en
población de América Latina. Por último, los mejores sistemas son ensambles grandes, los dos que generan texto no lo
comparan con reportes clínicos de referencia y solo \textcite{khan2025ensemble} marca los casos discordantes, sin
medir esa abstención.

\FloatBarrier
\subsection{Modelos visión--lenguaje en dermatología}
\label{sec:estado_arte_vlm}

Los \glspl{vlm} llegan a la dermatología por las dos vías de la Sección~\ref{sec:marco_vlm}: la alineación
contrastiva, que permite clasificar sin ajuste y recuperar casos, y la unión de un codificador visual con un modelo
de lenguaje que genera texto. Interesa de cada trabajo qué problema ataca, con qué método, contra qué se comparó y si
hubo expertos. Se presentan en orden cronológico dentro de cada estrategia, y las
Tablas~\ref{tab:marco_sota_vlm} y~\ref{tab:marco_sota_vlm_gen} los resumen.

\begin{table}[!t]
    \centering
    \caption{\glspl{vlm} contrastivos y modelos fundacionales en dermatología. $^\dagger$: sin revisión de pares.}
    \label{tab:marco_sota_vlm}
    \footnotesize
    \setlength{\tabcolsep}{3.5pt}
    \begin{tabularx}{\textwidth}{@{}>{\raggedright\arraybackslash}p{1.9cm}
        c
        >{\raggedright\arraybackslash\hspace{0pt}}X >{\raggedright\arraybackslash\hspace{0pt}}X
        >{\raggedright\arraybackslash\hspace{0pt}}p{2.2cm} >{\raggedright\arraybackslash\hspace{0pt}}X
        >{\raggedright\arraybackslash\hspace{0pt}}p{1.9cm}@{}}
        \toprule
        \textbf{Trabajo} & \textbf{Año} & \textbf{Arquitectura} & \textbf{Datos} & \textbf{Métricas} & \textbf{Resultado principal} & \textbf{Expertos} \\
        \midrule
        \textcite{kim2024transparent} & 2024 & CLIP ajustado de forma contrastiva
            & 105\,550 imágenes con texto de PubMed y libros
            & \gls{auc} por concepto
            & Tono de piel: 0.805 contra 0.736 (CLIP)
            & 2 dermatólogos verifican la recuperación \\
        \filasep
        \textcite{yan2025derm1miccv} & 2025 & DermLIP, tipo CLIP, cuatro codificadores
            & Derm1M, ontología con base de 4 dermatólogos
            & Exactitud, \gls{auc}, \emph{Recall@K}
            & \emph{Zero-shot} 0.561 contra 0.441 (BiomedCLIP)
            & Solo la ontología \\
        \filasep
        \textcite{yan2025makemiccai} & 2025 & CLIP con subtextos, alineación local y ponderación
            & 403\,563 pares imagen--texto
            & Exactitud, \gls{auc}, \emph{Recall@K}
            & \emph{Zero-shot} 0.493 contra 0.442 (CLIP, mismo corpus)
            & Ninguno \\
        \filasep
        \textcite{yan2025vision} & 2025 & ViT-Large auto-supervisado con maestro CLIP; no generativo
            & Más de 2 millones de imágenes, 11 instituciones
            & \gls{auc}, F1 ponderado
            & Mejora en 9 de 10 conjuntos; HAM10000 por 4.7\,\%
            & 3 estudios de lectura \\
        \filasep
        \textcite{yan2026dermfmzero}$^\dagger$ & 2026 & ViT-Large auto-supervisado y contrastivo; no generativo
            & Más de 3 millones de imágenes; 1\,014\,085 pares imagen--texto
            & F1 ponderado, \gls{auc}, sensibilidad macro, \emph{Recall@K}
            & \emph{Zero-shot} 0.733 contra 0.561 (DermLIP); atención primaria con ayuda, de 0.266 a 0.482
            & Tres estudios, con 30, 1\,090 y 34 clínicos \\
        \bottomrule
    \end{tabularx}
\end{table}

\begin{table}[!t]
    \centering
    \caption{\glspl{vlm} generativos en dermatología y su evaluación crítica. $^\dagger$: sin revisión de pares. ID: dentro de la distribución; OOD: fuera de ella.}
    \label{tab:marco_sota_vlm_gen}
    \footnotesize
    \setlength{\tabcolsep}{3.5pt}
    \begin{tabularx}{\textwidth}{@{}>{\raggedright\arraybackslash}p{1.9cm}
        c
        >{\raggedright\arraybackslash\hspace{0pt}}X >{\raggedright\arraybackslash\hspace{0pt}}X
        >{\raggedright\arraybackslash\hspace{0pt}}p{2.2cm} >{\raggedright\arraybackslash\hspace{0pt}}X
        >{\raggedright\arraybackslash\hspace{0pt}}p{1.9cm}@{}}
        \toprule
        \textbf{Trabajo} & \textbf{Año} & \textbf{Arquitectura} & \textbf{Datos} & \textbf{Métricas} & \textbf{Resultado principal} & \textbf{Expertos} \\
        \midrule
        \textcite{li2023llavamed} & 2023 & LLaVA de 7B ajustado en dos etapas
            & 600\,000 pares de PubMed Central; instrucciones de GPT-4
            & Puntaje relativo con GPT-4, exactitud en VQA
            & Supera al supervisado en VQA-RAD y PathVQA cerradas; sin evaluar en piel
            & Ninguno \\
        \filasep
        \textcite{zeng2025mmskin} & 2025 & LLaVA-Med de 7B con LoRA
            & 171\,600 imágenes de libros y nueve conjuntos públicos
            & BLEU-4, BERTScore, exactitud
            & \emph{Zero-shot} Patch16 82.34\,\% contra 68.55\,\% (LLaVA-Med)
            & Ninguno \\
        \filasep
        \textcite{shen2025skingptr1}$^\dagger$ & 2025 & Vision-R1-7B congelado con adaptadores
            & 334\,618 muestras con razonamiento sintético
            & Exactitud, puntajes de 0 a 5, peor grupo
            & 82.50\,\% en Derm12345; media de 3.6 sobre 5
            & 5 dermatólogos, 1\,000 casos \\
        \filasep
        \textcite{ru2026dermogpt}$^\dagger$ & 2026 & Qwen3-VL-8B con LoRA y aprendizaje por refuerzo
            & 772\,675 instrucciones; 33\,999 preguntas
            & Exactitud, juez automático, equidad
            & ID 78.04\,\% y OOD 66.48\,\% (Gemini-2.5-Flash: 63.37\,\% y 61.21\,\%)
            & 2 dermatólogos revisan; referencia humana en 100 casos por tarea \\
        \filasep
        \textcite{jiang2026multimodal}$^\dagger$ & 2026 & Evaluación sin ajuste de cinco modelos
            & Tres conjuntos públicos y 5\,811 consultas de un hospital
            & Exactitud de los tres primeros diagnósticos, triaje
            & Imagen sola en consulta real: 1.50\,\% a 13.35\,\% (abiertos), 24.65\,\% (GPT-4.1)
            & 3 dermatólogos validan el juez \\
        \bottomrule
    \end{tabularx}
\end{table}

\textcite{kim2024transparent} proponen MONET para anotar conceptos clínicos (eritema, úlcera, pigmentación) sobre
imágenes sin etiquetas manuales, con el fin de auditar conjuntos de datos y modelos. Es una versión de CLIP ajustada
de forma contrastiva con 105\,550 imágenes dermatológicas y sus descripciones, extraídas de PubMed y de libros de
texto. Se compara con CLIP y con un ResNet-50 supervisado con conceptos anotados, sobre SkinCon y Derm7pt. En la
anotación del tono de piel alcanza un \gls{auc} de 0.805 contra 0.736 de CLIP y 0.874 del supervisado, que se evalúa
en condiciones más favorables. Dos dermatólogos verifican la recuperación de imágenes por concepto, y un clasificador
por conceptos sobre MONET supera, en imágenes clínicas, a las líneas de base en malignidad (\gls{auc} 0.815) y melanoma (0.881). Los autores
aclaran que no sirve para diagnosticar, que fallará con conceptos ausentes del entrenamiento y que no pudieron medir
el tono de piel en imágenes dermatoscópicas.

\textcite{yan2025derm1miccv} entrenan la familia DermLIP, modelos de tipo CLIP con la pérdida contrastiva original, sobre
Derm1M, un corpus de alrededor de un millón de pares imagen--texto dermatológicos organizados según una ontología
cuya base crearon cuatro dermatólogos y que se amplió con ayuda de un modelo de lenguaje. El trabajo, presentado en ICCV 2025, compara con CLIP, SigLIP,
CoCa, PMC-CLIP, BiomedCLIP y MONET en ocho conjuntos y cuatro tareas. En clasificación \emph{zero-shot}, la
exactitud promedio sobre cuatro conjuntos es 0.561 contra 0.441 de BiomedCLIP, y 0.588 con un codificador visual
dermatológico (PanDerm). En la anotación \emph{zero-shot} de conceptos supera a MONET en 1.1 puntos de \gls{auc} en promedio, con
un codificador visual más chico. Los autores reconocen el ruido que introducen los textos educativos, y no hubo evaluación con
clínicos.

\textcite{yan2025makemiccai} observan que CLIP trunca los textos a 77 \emph{tokens} y descarta detalle clínico. MAKE
divide cada descripción en oraciones, a las que suma un texto de enfermedad y otro de conceptos generados por un
modelo de lenguaje; alinea esos subtextos con regiones de la imagen y pondera cada uno según su relevancia
diagnóstica. Se preentrena con 403\,563 pares y se compara con CLIP, SigLIP y CoCa entrenados con el mismo
corpus y con BiomedCLIP, PMC-CLIP y MONET. La exactitud \emph{zero-shot} promedio en cinco conjuntos sube de 0.442
(CLIP sobre el mismo corpus) a 0.493, el \gls{auc} de conceptos mejora 1.11 puntos sobre la mejor línea de base y en
recuperación llega a 0.444 de \emph{Recall@50}. No hay evaluación con expertos ni discusión de limitaciones.

\textcite{yan2025vision} proponen PanDerm, un modelo fundacional visual y no generativo. Se preentrena sin etiquetas
con más de dos millones de imágenes de 11 instituciones y cuatro modalidades, con modelado latente enmascarado y
alineación con las características de CLIP como maestro, sobre cuatro GPU H100. Se evalúa en 28 tareas, contra
DINOv2, SwAVDerm, CLIP, MONET, BiomedCLIP y BiomedGPT. Mejora en 9 de 10 conjuntos de diagnóstico (por 4.7\,\% en
HAM10000 y 9.0\,\% en PAD-UFES-20 sobre el siguiente modelo), y con solo el 10\,\% de las etiquetas suele
superar a los demás. Tres estudios de lectura con clínicos muestran que, con imágenes secuenciales de melanoma, su exactitud supera en
10.2 puntos a la del revisor promedio entre 12 clínicos, que con su ayuda la exactitud con dermatoscopía sube 11
puntos y que la exactitud de los tres primeros diagnósticos de profesionales no dermatólogos sube 16.5 puntos en 128
afecciones. Los autores reconocen que cubren unas 200 de más de 1\,000
afecciones y que la evaluación de equidad es incompleta.

\textcite{yan2026dermfmzero}, disponible en arXiv y Research Square, sin revisión de pares, del mismo grupo, busca que un modelo
fundacional visión--lenguaje apoye la decisión clínica sin ajuste por tarea y evaluarlo con clínicos. Es un
ViT-Large no generativo, preentrenado de forma auto-supervisada con más de tres millones de imágenes y después de
forma contrastiva con 1\,014\,085 pares imagen--texto. Agrega autoencoders esparsos que nombran conceptos y permiten
suprimir las neuronas asociadas a artefactos como reglas y pelo. En siete conjuntos \emph{zero-shot} el promedio de
\gls{auc} (conjuntos binarios) y exactitud balanceada (multiclase) es 0.733 contra 0.561 de DermLIP, y en HAM10000 la sensibilidad macro es 0.744 contra 0.512. En tres estudios con
clínicos, 30 médicos de atención primaria con su ayuda pasan de 0.266 a 0.482 de exactitud en los tres primeros
diagnósticos. Los autores reconocen que los estudios son retrospectivos y simulados y que el preentrenamiento cubre
más de 400 de las más de 2\,000 entidades dermatológicas.

La segunda estrategia une un codificador visual con un modelo de lenguaje. \textcite{li2023llavamed} presentan
LLaVA-Med, un asistente biomédico de propósito general. Parte de LLaVA de 7\,000 millones de parámetros y lo ajusta
en dos etapas con figuras de PubMed Central: primero con 600\,000 pares figura--texto, para alinear el vocabulario,
y después con 60\,000 conversaciones generadas por GPT-4 a partir de los textos, en menos de 15 horas con ocho A100. Se evalúa en tres
conjuntos de preguntas visuales de radiología y patología y con un puntaje relativo calculado por GPT-4. Tras el ajuste
supera al estado del arte supervisado en las preguntas cerradas de VQA-RAD y PathVQA, y en las abiertas solo en SLAKE.
No se evalúa en dermatología ni con expertos, y sus autores reconocen alucinaciones y poco razonamiento en profundidad.
Med-Flamingo, PeFoMed y LLaVA-Tri también se evalúan en preguntas visuales de radiología y patología, y
Med-Flamingo además en problemas clínicos de varias especialidades, pero ninguno en
piel~\cite{moor2023medflamingo,he2025pefomed,xie2025medtrinity25m}.

\textcite{zeng2025mmskin} adaptan ese modelo a la piel con SkinVL. Sobre LLaVA-Med de 7\,000 millones de parámetros
entrenan solo la proyección visual y adaptadores LoRA, con un corpus de imágenes de libros de texto y nueve conjuntos
públicos de diagnóstico reformulados como preguntas (171\,600 imágenes en la variante completa). Se compara con cinco
modelos generales y dos modelos médicos de propósito general, entre ellos LLaVA-Med y HuatuoGPT, en descripción, preguntas y
clasificación con ajuste y \emph{zero-shot}. En \emph{zero-shot} la variante completa alcanza 82.34\,\% de exactitud en
Patch16 contra 68.55\,\% de LLaVA-Med, y es el único modelo que supera 50\,\% en HAM10000. Los textos provienen de
libros de dermatólogos, pero no hubo evaluación de las respuestas por expertos ni se discuten limitaciones, más allá de
pedir mayor capacidad de razonamiento.

\textcite{shen2025skingptr1}, publicado solo en arXiv y sin revisión de pares, buscan un modelo que razone de forma
explícita, produzca un reporte estructurado y rinda de forma pareja entre tonos de piel. Mantienen congelado
Vision-R1-7B y entrenan solo el 0.056\,\% de los parámetros, con un sesgo sobre los \emph{logits}, un adaptador de
mezcla de expertos que incorpora un tono de piel estimado y la destilación de PanDerm. El entrenamiento supervisado usa
334\,618 muestras con razonamiento sintetizado por modelos de lenguaje. En siete conjuntos externos, contra cinco
sistemas entre ellos GPT-4o mini y LLaVA-Med, alcanza 82.50\,\% de exactitud en Derm12345 y entre 42.40\,\% y
69.30\,\% en otros tres. Cinco dermatólogos califican 1\,000 casos con una media de 3.6 sobre 5, pero solo sobre las
salidas de este modelo. Los autores reconocen bucles de autovalidación con incertidumbre alta y una latencia que impide
ejecutarlo en dispositivos de borde.

\textcite{ru2026dermogpt}, publicado solo en arXiv y sin revisión de pares, proponen un modelo que siga el orden
clínico de morfología, razonamiento y diagnóstico, y anuncian que publicarán los datos, el \emph{benchmark} y los
pesos tras la aceptación. Reúnen DermoInstruct,
con 772\,675 instrucciones, y DermoBench, con 33\,999 preguntas. Ajustan Qwen3-VL-8B con LoRA y después con aprendizaje
por refuerzo, con una recompensa que combina la exactitud, la similitud jerárquica dentro de la ontología diagnóstica y
la coincidencia de atributos morfológicos. Comparan 16 modelos. En diagnóstico alcanza 78.04\,\% dentro de la
distribución y 66.48\,\% fuera de ella, contra 63.37\,\% y 61.21\,\% de Gemini-2.5-Flash, el mejor modelo comercial.
Dos dermatólogos revisan línea por línea un núcleo de 900 casos, y dos clínicos de dermatología resuelven 100 casos
por tarea como referencia humana. Los autores reconocen que la
brecha con la referencia humana persiste.

La tercera línea mide qué tan lejos están estos modelos de la consulta real. \textcite{jiang2026multimodal}, también
publicado solo en arXiv y sin revisión de pares, y de un grupo independiente de los anteriores, evalúan sin ajuste
cinco modelos (LLaVA-Med v1.5, SkinGPT-4, InternVL-Chat, MedGemma-4B y GPT-4.1) en diagnóstico diferencial y triaje por
gravedad. Usan tres conjuntos públicos y una cohorte de 5\,811 consultas de un hospital, con 46\,405 imágenes, que no
se distribuye. La exactitud de los tres primeros diagnósticos la calcula un juez automático validado contra tres
dermatólogos. En las consultas con la imagen sola los modelos abiertos obtienen entre 1.50\,\% y 13.35\,\% y GPT-4.1
24.65\,\%. Sin contexto, la sensibilidad de triaje queda por debajo de 0.50 en todos, y el texto del clínico remitente
mejora los resultados pero los vuelve sensibles a errores de ese texto. Los autores concluyen que no corresponde
usarlos sin supervisión y reconocen que los casos provienen de un servicio con internación y que son solo cinco
modelos.

La línea contrastiva avanzó de adaptar CLIP con texto de la literatura (MONET) a corpus dermatológicos con ontología
(DermLIP), textos descompuestos por aspecto (MAKE) y preentrenamiento visual a gran escala (PanDerm, DermFM-Zero), y
solo estos dos últimos se evalúan en estudios de lectura con clínicos. La línea generativa pasó de asistentes biomédicos que no se evaluaron en
piel (LLaVA-Med) a modelos adaptados a ella (SkinVL) y a modelos que razonan y producen un reporte estructurado
(SkinGPT-R1, DermoGPT). La evaluación crítica muestra que los \glspl{vlm} médicos de propósito general rinden poco en
piel y que la brecha crece en consultas reales. Quedan cuestiones abiertas.
Ningún trabajo combina un clasificador supervisado con un generador ni recupera casos parecidos para anclar el texto.
La incertidumbre aparece como limitación, o en DermoGPT para ponderar sus predicciones, y no como una decisión
entre afirmar, ofrecer diferenciales o derivar (Sección~\ref{sec:marco_tres_salidas}). El cómputo es alto, de ocho RTX 4090 para ajustar SkinGPT-R1 a cuatro H100 para
preentrenar PanDerm o DermFM-Zero. Y los estudios con clínicos son pocos: solo PanDerm y DermFM-Zero miden el efecto de
la ayuda sobre sus decisiones, y en los generativos los expertos son dos o tres, o evalúan solo las salidas del propio modelo.

\FloatBarrier
\subsection{Generación de reportes dermatológicos}
\label{sec:estado_arte_reportes}

Un reporte dermatológico reúne la descripción de lo que se ve, un diagnóstico y, cuando corresponde, una
recomendación. Los trabajos difieren en cómo lo producen (un único \gls{vlm} o un clasificador que precede a un
redactor), en la evidencia que usan (solo la imagen o también un contexto recuperado) y en cómo evalúan el texto
(Sección~\ref{sec:marco_rep_ejes}). La Tabla~\ref{tab:marco_generadores} resume los generadores, y la
Tabla~\ref{tab:marco_sota_reportes}, los trabajos que se ocupan de la evaluación y de la recuperación.

\begin{table}[!t]
    \centering
    \caption{Generadores de texto dermatológico. «Métricas con referencia»: BLEU, ROUGE y BERTScore; «---»: no
    informado.}
    \label{tab:marco_generadores}
    \small
    \begin{tabularx}{\textwidth}{@{}>{\raggedright\arraybackslash}p{3.2cm}
        >{\raggedright\arraybackslash\hspace{0pt}}X >{\raggedright\arraybackslash\hspace{0pt}}X l@{}}
        \toprule
        \textbf{Trabajo} & \textbf{Modelo que redacta} & \textbf{Evaluación del texto} & \textbf{Hardware} \\
        \midrule
        SkinGPT-4~\cite{zhou2024pretrained} & MiniGPT-4 con Llama-2-13B & dermatólogos sobre 150 casos reales & 8 A100 \\
        \filasep
        SkinVL~\cite{zeng2025mmskin} & LLaVA-Med con LoRA & métricas con referencia & 8 RTX 3090 \\
        \filasep
        DermatoLlama~\cite{yilmaz2025resourceefficient} & Llama 3.2 11B con LoRA & exactitud y métricas con referencia & 4 RTX 4090 \\
        \filasep
        Skin-R1~\cite{liu2025skinr1} & Qwen2.5-VL-7B & exactitud del diagnóstico en opción múltiple & 2 A100 \\
        \filasep
        SkinGPT-R1~\cite{shen2025skingptr1} & Vision-R1-7B congelado con adaptadores & puntajes de un juez automático y de 5 dermatólogos & 8 RTX 4090 \\
        \filasep
        Khan et al.~\cite{khan2025ensemble} & Llama-3 70B por API & --- & --- \\
        \filasep
        SkinSavvy2~\cite{kim2025skinsavvy2} & GPT-4 & encuesta a 20 participantes, 2 dermatólogos & 4 RTX A6000 \\
        \filasep
        Chakrabarti et al.~\cite{chakrabarti2025skin} & LLaVA-1.5-13B destilado de GPT-4V & métricas con referencia, corrección diagnóstica y revisión de un médico & --- \\
        \filasep
        MpoxVLM~\cite{cao2024mpoxvlm} & Llama-2-7B con LoRA & exactitud del diagnóstico & 2 RTX 4090 \\
        \bottomrule
    \end{tabularx}
\end{table}

\begin{table}[!t]
    \centering
    \caption{Evaluación de narrativas dermatológicas y recuperación de contexto para \glspl{vlm} médicos.}
    \label{tab:marco_sota_reportes}
    \footnotesize
    \setlength{\tabcolsep}{3.5pt}
    \begin{tabularx}{\textwidth}{@{}>{\raggedright\arraybackslash}p{1.9cm}
        c
        >{\raggedright\arraybackslash\hspace{0pt}}X >{\raggedright\arraybackslash\hspace{0pt}}X
        >{\raggedright\arraybackslash\hspace{0pt}}p{2.2cm} >{\raggedright\arraybackslash\hspace{0pt}}X
        >{\raggedright\arraybackslash\hspace{0pt}}p{1.9cm}@{}}
        \toprule
        \textbf{Trabajo} & \textbf{Año} & \textbf{Arquitectura} & \textbf{Datos} & \textbf{Métricas} & \textbf{Resultado principal} & \textbf{Expertos} \\
        \midrule
        \textcite{shen2025dermbench} & 2025 & Juez automático y evaluador LLaVA v1.5 7B con LoRA
            & 4\,000 imágenes de DermNet con referencia; 500 imágenes y nueve modelos
            & Error absoluto medio contra médicos, seis dimensiones
            & Error medio 0.363 (DermBench) y 0.178 (DermEval)
            & 2 dermatólogos, coautores \\
        \filasep
        \textcite{xia2025mmedrag} & 2025 & LLaVA-Med-1.5 con recuperación de reportes y optimización de preferencias
            & Radiología, oftalmología y patología
            & Exactitud, F1, \gls{auc}; BLEU, ROUGE-L, METEOR
            & IU-Xray: exactitud 75.47 a 89.54; BLEU 9.64 a 31.38
            & No \\
        \filasep
        \textcite{xia2024rule} & 2024 & LLaVA-Med-1.5 con recuperación de reportes, $k$ calibrado y optimización
            de preferencias
            & Radiología y oftalmología
            & Exactitud, F1; BLEU, ROUGE-L, METEOR
            & IU-Xray: exactitud 75.47 a 87.84; BLEU 9.64 a 27.53
            & No \\
        \bottomrule
    \end{tabularx}
\end{table}

Un primer grupo genera el texto con un único \gls{vlm} ajustado a dermatología. SkinVL y SkinGPT-R1 pertenecen a
este grupo y se describieron en la Sección~\ref{sec:estado_arte_vlm}; aquí se detallan otros tres.

SkinGPT-4 busca que una persona sin acceso a un dermatólogo suba una foto y obtenga en una conversación la
descripción de la lesión, un diagnóstico y recomendaciones. Ajusta MiniGPT-4 con Llama-2-13B en dos pasos: primero
alinea la imagen con conceptos clínicos y después aprende diagnósticos a partir de notas de médicos, con unas 50\,000
imágenes públicas y propias~\cite{zhou2024pretrained}. Lo compara con MiniGPT-4 sin adaptar y con
cada paso por separado, en ejemplos cualitativos. El texto lo evalúan cinco dermatólogos certificados con un cuestionario
sobre 150 casos reales, sin métricas con referencia: consideran correcto o pertinente el 80.63\,\% de los
diagnósticos. El entrenamiento llevó unas 24 horas en 8 A100 de 80\,GB. Los autores reconocen que los tonos
Fitzpatrick V y VI están poco representados, que las enfermedades con varias lesiones combinadas siguen siendo
difíciles y que las alucinaciones de los modelos de lenguaje pueden dar información errónea al paciente.

DermatoLlama apunta a un modelo que se ejecute localmente, por privacidad y por costo, y lo logra con datos
sintéticos~\cite{yilmaz2025resourceefficient}. A partir de 367 tareas semilla revisadas por expertos y 82\,379
imágenes abiertas, modelos más grandes consultados por API generan 1.2 millones de textos sintéticos (DermaSynth),
con los que ajusta con LoRA un Llama~3.2 de 11 mil millones de parámetros. Se compara con el modelo base, GPT-4o y
Gemini~2.0~Flash en tres evaluaciones. La primera, de 225 imágenes de DERM12345, la anotaron tres dermatólogos
coautores con preguntas jerárquicas; allí DermatoLlama alcanza una exactitud de 0.87 (GPT-4o, 0.86), un BLEU-4 de
0.66 (0.55) y un ROUGE-L de 0.93 (0.81). Se mide, pues, el diagnóstico y el parecido con un texto de referencia, y
ningún experto revisa los textos generados. Entrena en 4 RTX 4090 y ejecuta en una. Los autores advierten que los
datos sintéticos pueden perder rasgos sutiles en casos raros y que faltan la validación clínica externa y la
auditoría de sesgos. Está publicado solo en medRxiv, sin revisión de pares.

Skin-R1 se ocupa de que el diagnóstico venga con un razonamiento fundado, no de redactar un reporte
completo~\cite{liu2025skinr1}. Sobre Qwen2.5-VL-7B con LoRA, genera trayectorias de razonamiento a partir de libros
de texto, con la jerarquía de enfermedades y el diagnóstico diferencial; las usa para el ajuste supervisado y luego
para un aprendizaje por refuerzo con recompensas por exactitud jerárquica, por discriminación de malignidad y
por formato. Lo evalúa con preguntas de
opción múltiple sobre seis conjuntos de datos de dermatología y un conjunto externo de dermatoscopía. Se compara con
\glspl{vlm} generales (Qwen, InternVL, LLaVA), médicos (LLaVA-Med, HuatuoGPT-Vision, MedGemma) y SkinVL, que casi
nunca produce respuestas válidas con ese formato. La exactitud media es 0.6385 en las pruebas de los mismos seis conjuntos
del entrenamiento (0.1955 sobre el siguiente modelo) y 0.7171 en el conjunto externo; en la clasificación benigno, maligno o precanceroso, 0.6928 de
exactitud y 0.4287 de F1 macro. Entrena en 2 A100 de 40\,GB. No mide la calidad del texto del razonamiento ni
incluye expertos, y está publicado solo en arXiv, sin revisión de pares.

Un segundo grupo separa el diagnóstico de la redacción: un clasificador decide y un modelo de lenguaje escribe.
Los dos sistemas de este grupo descriptos en la Sección~\ref{sec:estado_arte_diagnostico} redactan con un
modelo grande. En el de \textcite{khan2025ensemble}, Llama-3 de 70 mil millones de parámetros, consultado por API,
recibe solo el nombre de la enfermedad, el tipo de consenso entre los clasificadores y la confianza, no la imagen, y
redacta la descripción de la lesión, los síntomas por vigilar, el tratamiento y las señales de alarma; el texto solo
se valora con una rúbrica propia de contenido, sin métricas con referencia ni expertos. En SkinSavvy2, la clase predicha se combina en un \emph{prompt} con la edad,
el sexo, el tipo de piel, la exposición solar y los antecedentes familiares, y GPT-4 redacta una recomendación de
cuidado~\cite{kim2025skinsavvy2}. Veinte participantes, 12 legos, 6 profesionales de la salud y 2 dermatólogos,
puntuaron de 0 a 5 la confianza, la personalización, la utilidad y la guía de manejo de 20 respuestas, y el
\emph{prompt} con datos personales y GPT-4o obtuvo los puntajes más altos. Los dermatólogos lo consideraron una ayuda
y no un reemplazo del diagnóstico profesional, y los autores reconocen que los datos que ingresa el usuario pueden ser
inexactos.

Chakrabarti et al.\ generan reportes estructurados (hallazgos, impresión y diagnóstico diferencial) de imágenes
dermatoscópicas según la lista de 7 puntos para melanoma~\cite{chakrabarti2025skin}. Sin pares imagen-reporte
disponibles, GPT-4V actúa como maestro: a partir de la imagen, una descripción de libro de texto, los rasgos anotados
de la lista y el diagnóstico, produce pseudorreportes con los que se ajusta LLaVA-1.5-13B, sobre 305 imágenes. Se
compara con Claude~3.7 Sonnet, con LLaVA sin ajustar y con dos variantes de recuperación (BioMedCLIP y PubMedCLIP)
que anteponen al \emph{prompt} los pseudorreportes más parecidos. Contra los pseudorreportes de referencia, el modelo
obtiene SacreBLEU 55.59, ROUGE-1 0.5438 y BERTScore F1 0.9025; la recuperación con BioMedCLIP, 21.84, 0.5162 y
0.8966. Evalúa además los rasgos de la lista y la sensibilidad y especificidad para melanoma, y un médico revisa 30
reportes de este modelo y de GPT-4V. Como la referencia la escribe el propio maestro, las métricas miden el parecido
con GPT-4V. Los autores atribuyen el SacreBLEU inusualmente alto a la estructura de plantilla de los
pseudorreportes y señalan que las métricas léxicas no bastan. Reconocen también el desbalance de clases y las confusiones entre lesiones pigmentadas benignas y melanoma.

SkinGEN acompaña el diagnóstico con imágenes de referencia generadas por difusión~\cite{lin2025skingen}, y MpoxVLM
combina la imagen con datos clínicos para diagnosticar una sola enfermedad, la viruela
símica~\cite{cao2024mpoxvlm}.

Dos líneas complementan a los generadores: cómo evaluar el texto y cómo condicionar la generación con contexto
recuperado.

\textcite{shen2025dermbench} abordan cómo evaluar las narrativas diagnósticas más allá de las métricas de similitud
textual, que según los autores correlacionan poco con la utilidad clínica. DermBench tiene una narrativa de
referencia para cada una de 4\,000 imágenes de DermNet, redactada por modelos de lenguaje a partir de la etiqueta
verdadera y corregida por dermatólogos. Un juez automático que no ve la imagen compara el candidato con la
referencia y lo puntúa de 0 a 5 en seis dimensiones. DermEval es un evaluador sin referencia, un LLaVA v1.5 7B con
LoRA entrenado con aprendizaje por refuerzo contra los puntajes de médicos. Sobre nueve modelos generales o médicos,
el error absoluto medio contra los médicos, promediado en las seis dimensiones, es 0.363 para DermBench y 0.178
para DermEval. Dos dermatólogos,
coautores, puntuaron todo, sin cuantificar el acuerdo entre ellos. Los autores reconocen que las referencias se
redactan con modelos, lo que puede introducir sesgo. El trabajo se presentó en ECCV 2026.

MMed-RAG~\cite{xia2025mmedrag} y RULE~\cite{xia2024rule}, del mismo grupo, son los antecedentes de la recuperación
para \glspl{vlm} médicos. Buscan mejorar la factualidad de LLaVA-Med-1.5 con contexto recuperado y evitar que el
modelo ignore la imagen o que un contexto equivocado degrade respuestas correctas. En los dos la consulta es la
imagen y lo recuperado son reportes de texto, que se concatenan en el \emph{prompt}; ninguno recupera imágenes ni se
abstiene. RULE calibra cuántos reportes recuperar con un control estadístico del riesgo, y MMed-RAG corta la lista
donde la similitud cae de forma brusca. Los dos ajustan además el generador con optimización directa de preferencias,
para que use el contexto solo cuando ayuda. Los datos son de radiología, oftalmología y, en MMed-RAG, patología; no
hay dermatología. Se comparan con métodos de decodificación y con otros \glspl{vlm} médicos, y MMed-RAG también
con métodos de recuperación, entre ellos RULE. En IU-Xray la exactitud
de las preguntas sube de 75.47 a 87.84 con RULE y a 89.54 con MMed-RAG. En la ablación de RULE, sin embargo, el
contexto recuperado sin más empeora MIMIC-CXR, de 75.79 a 67.35. Ninguno incluye expertos, y MMed-RAG señala que una
modalidad nueva puede exigir reentrenar el recuperador.

El texto se evalúa de tres maneras que casi no se combinan (Sección~\ref{sec:marco_eval_texto}). Las métricas con
referencia miden el parecido con un texto y no la corrección clínica, y las referencias suelen ser sintéticas:
DermatoLlama usa preguntas anotadas por sus coautores y Chakrabarti et al.\ los pseudorreportes del propio maestro.
La corrección del diagnóstico (Skin-R1, DermatoLlama) deja fuera la descripción y la recomendación. Los expertos
son pocos: cinco dermatólogos sobre 150 casos en SkinGPT-4, dos dermatólogos entre 20 participantes en SkinSavvy2,
un médico sobre 30 reportes en Chakrabarti et al.\ y dos coautores en DermBench; Khan et al.\ solo aplican una
rúbrica propia.
La recuperación puede mejorar la factualidad y también degradarla: en la ablación de RULE el contexto sin más
pierde 8.4 puntos en MIMIC-CXR, y en dermatología la recuperación de pseudorreportes rinde menos que el ajuste. En todos los casos
se recuperan textos y no imágenes, y las métricas miden el parecido con una referencia. El hardware va de dos GPU de
consumo a ocho A100 (DermatoLlama se ejecuta en una sola), y varios sistemas dependen de modelos cerrados o consultados por API
para redactar o para generar sus datos o referencias (Llama-3 70B en Khan et al., GPT-4 en SkinSavvy2, GPT-4V en
Chakrabarti et al., Gemini~2.0~Flash en DermatoLlama y Gemini entre otros modelos en DermBench).

\FloatBarrier
\subsection{Posicionamiento del trabajo}
\label{sec:estado_arte_posicion}

Cada línea deja algo sin resolver. Los clasificadores de lesiones descansan sobre la prueba de un desafío, un solo
centro o una partición temporal, y los de mejor desempeño son ensambles grandes. Entre los \glspl{vlm}, los
contrastivos especializados clasifican y recuperan pero no redactan, los generativos exigen un cómputo elevado y los
\glspl{vlm} médicos de propósito general rinden poco en piel. En los reportes, la mayoría de los sistemas se evalúa por parecido textual o por
la exactitud del diagnóstico, y la recuperación estudiada trae textos de otras especialidades.

De esos huecos sale el \emph{pipeline}: un clasificador que estima el diagnóstico, una recuperación de casos
parecidos y un generador de reportes, con una decisión que depende de la confianza y determina cuánto afirma el
texto. Cada pieza atiende una brecha de la Tabla~\ref{tab:marco_brechas}
(Sección~\ref{sec:marco_tres_salidas}), y el análisis incluye cuánto degrada un vecino equivocado el
texto, como advierten MMed-RAG y RULE~\cite{xia2025mmedrag,xia2024rule}.

\begin{table}[!t]
    \centering
    \caption{De las brechas de la revisión a las decisiones del trabajo.}
    \label{tab:marco_brechas}
    \small
    \begin{tabularx}{\textwidth}{@{}>{\raggedright\arraybackslash}p{3.3cm}
        >{\raggedright\arraybackslash\hspace{0pt}}X >{\raggedright\arraybackslash\hspace{0pt}}X@{}}
        \toprule
        \textbf{Brecha} & \textbf{Evidencia en la revisión} & \textbf{Decisión del trabajo} \\
        \midrule
        Sin comparación bajo un mismo protocolo
            & Clasificadores casi solo convolucionales, validados sobre pruebas de desafíos o en un centro
              (Sección~\ref{sec:estado_arte_diagnostico}); \glspl{vlm} medidos entre sí o con clínicos
              (Sección~\ref{sec:estado_arte_vlm})
            & Comparar clasificadores supervisados y modelos médicos generales y dermatológicos con un mismo
              protocolo, con partición disjunta por lesión y un conjunto externo \\
        \filasep
        Texto evaluado por su parecido con una referencia
            & Métricas con referencia, exactitud del diagnóstico o un juez que no ve la imagen
              (Sección~\ref{sec:estado_arte_reportes})
            & Medir por separado la fidelidad diagnóstica y la calidad descriptiva
              (Sección~\ref{sec:marco_rep_ejes}) \\
        \filasep
        Incertidumbre y evidencia sin integrar
            & La incertidumbre como limitación y no como decisión; contexto recuperado de otras especialidades
              (Secciones~\ref{sec:estado_arte_vlm} y~\ref{sec:estado_arte_reportes})
            & Recuperar casos parecidos por imagen, usar la concordancia como señal de confianza y decidir entre
              afirmar, ofrecer diferenciales o derivar
              (Secciones~\ref{sec:marco_integracion} y~\ref{sec:marco_tres_salidas}) \\
        \filasep
        Cómputo fuera del alcance de un equipo sin infraestructura
            & Ocho RTX 4090, cuatro H100, ensambles de 90 redes o un modelo de 70 mil millones de parámetros
              por API
            & Ejecutar todo en una GPU de consumo, con un codificador contrastivo dermatológico y LoRA \\
        \bottomrule
    \end{tabularx}
\end{table}

El codificador contrastivo es DermLIP. DermFM-Zero lo supera en su propio protocolo, pero DermLIP, en su variante
ViT-B/16, se ejecuta en una GPU de consumo, y un mismo espacio sirve para clasificar sin ajuste, entrenar un
clasificador lineal y recuperar vecinos visuales (Sección~\ref{sec:marco_recuperacion_imagenes},
\cite{yan2025derm1miccv}). Los modelos generativos se adaptan con LoRA (Sección~\ref{sec:marco_llm_lora}).

Dos trabajos se acercan a una parte del diseño. El ensamble de \textcite{khan2025ensemble} deriva a un especialista
los casos en que sus miembros no coinciden, pero la concordancia solo agrega al texto un aviso de revisión, y el
texto lo escribe un modelo consultado por API. \textcite{chakrabarti2025skin} miden la corrección diagnóstica de los reportes de un único
modelo destilado. Ninguno combina en una cascada la recuperación de casos, la decisión por concordancia y la
redacción. La implementación y el protocolo de evaluación se describen en el Capítulo~\ref{cp:metodos}.

\FloatBarrier

